
\documentclass[11pt]{article}
\usepackage{amssymb}
\usepackage{amsmath}
\usepackage{amsfonts}
\usepackage{setspace}
\usepackage[margin=1.0in,nomarginpar,nohead]{geometry}
\usepackage{verbatim}
\usepackage[table]{xcolor}
\usepackage[longnamesfirst,authoryear]{natbib}

\setcounter{MaxMatrixCols}{10}

\definecolor{myEvenRowsColor}{rgb}{0.93,0.93,1}
\newtheorem{theorem}{Theorem}
\newtheorem{acknowledgement}[theorem]{Acknowledgement}
\newtheorem{algorithm}[theorem]{Algorithm}
\newtheorem{axiom}[theorem]{Axiom}
\newtheorem{case}[theorem]{Case}
\newtheorem{claim}[theorem]{Claim}
\newtheorem{conclusion}[theorem]{Conclusion}
\newtheorem{condition}[theorem]{Condition}
\newtheorem{conjecture}[theorem]{Conjecture}
\newtheorem{corollary}[theorem]{Corollary}
\newtheorem{criterion}[theorem]{Criterion}
\newtheorem{definition}[theorem]{Definition}
\newtheorem{example}[theorem]{Example}
\newtheorem{exercise}[theorem]{Exercise}
\newtheorem{lemma}[theorem]{Lemma}
\newtheorem{notation}[theorem]{Notation}
\newtheorem{problem}[theorem]{Problem}
\newtheorem{proposition}[theorem]{Proposition}
\newtheorem{remark}[theorem]{Remark}
\newtheorem{solution}[theorem]{Solution}
\newtheorem{summary}[theorem]{Summary}
\onehalfspacing

\RequirePackage{graphicx}
\RequirePackage[normalem]{ulem}
\RequirePackage{xcolor}
\RequirePackage{amssymb}
\RequirePackage{amsmath}
\RequirePackage{amsfonts}
\RequirePackage{geometry}
\RequirePackage{xspace} %
\RequirePackage{verbatim}
\RequirePackage{environ}
\RequirePackage{bbm}
\RequirePackage{hyperref}

\providecommand{\III}[1]{\mathbb{I}\left\{#1\right\}}%

\newcommand*{\II}{\mathbbm{1}} %

\newcommand\formatAppx{
\renewcommand{\theequation}{A.\arabic{equation}}
\renewcommand{\thelemma}{A.\arabic{lemma}}
\renewcommand{\thetheorem}{A.\arabic{theorem}}
\renewcommand{\thesection}{\Alph{section}}
\setcounter{equation}{0}%
\setcounter{theorem}{0} \setcounter{section}{0}
}

\newcommand\CMDsmall\small %

\makeatletter
                                     %
                                     %
                                     %
\renewcommand\subparagraph{\@startsection{subparagraph}{5}{\z@}%
                                     {-3.25ex\@plus -1ex \@minus -.2ex}%
                                     {0.0001pt \@plus .2ex}%
                                     {\normalfont\normalsize\bfseries}}
\makeatother

\hypersetup{colorlinks=true,anchorcolor=blue,citecolor=black,linkcolor=[rgb]{0,0.0,0.75}} %

\sloppy
% Macros for Scientific Word and Scientific WorkPlace 5.5 documents saved with the LaTeX filter.
% Copyright (C) 2005 Mackichan Software, Inc.

\typeout{TCILATEX Macros for Scientific Word and Scientific WorkPlace 5.5 <06 Oct 2005>.}
\typeout{NOTICE:  This macro file is NOT proprietary and may be 
freely copied and distributed.}
%
\makeatletter

%%%%%%%%%%%%%%%%%%%%%
% pdfTeX related.
\ifx\pdfoutput\relax\let\pdfoutput=\undefined\fi
\newcount\msipdfoutput
\ifx\pdfoutput\undefined
\else
 \ifcase\pdfoutput
 \else 
    \msipdfoutput=1
    \ifx\paperwidth\undefined
    \else
      \ifdim\paperheight=0pt\relax
      \else
        \pdfpageheight\paperheight
      \fi
      \ifdim\paperwidth=0pt\relax
      \else
        \pdfpagewidth\paperwidth
      \fi
    \fi
  \fi  
\fi

%%%%%%%%%%%%%%%%%%%%%
% FMTeXButton
% This is used for putting TeXButtons in the 
% frontmatter of a document. Add a line like
% \QTagDef{FMTeXButton}{101}{} to the filter 
% section of the cst being used. Also add a
% new section containing:
%     [f_101]
%     ALIAS=FMTexButton
%     TAG_TYPE=FIELD
%     TAG_LEADIN=TeX Button:
%
% It also works to put \defs in the preamble after 
% the \input tcilatex
\def\FMTeXButton#1{#1}
%
%%%%%%%%%%%%%%%%%%%%%%
% macros for time
\newcount\@hour\newcount\@minute\chardef\@x10\chardef\@xv60
\def\tcitime{
\def\@time{%
  \@minute\time\@hour\@minute\divide\@hour\@xv
  \ifnum\@hour<\@x 0\fi\the\@hour:%
  \multiply\@hour\@xv\advance\@minute-\@hour
  \ifnum\@minute<\@x 0\fi\the\@minute
  }}%

%%%%%%%%%%%%%%%%%%%%%%
% macro for hyperref and msihyperref
%\@ifundefined{hyperref}{\def\hyperref#1#2#3#4{#2\ref{#4}#3}}{}

\def\x@hyperref#1#2#3{%
   % Turn off various catcodes before reading parameter 4
   \catcode`\~ = 12
   \catcode`\$ = 12
   \catcode`\_ = 12
   \catcode`\# = 12
   \catcode`\& = 12
   \catcode`\% = 12
   \y@hyperref{#1}{#2}{#3}%
}

\def\y@hyperref#1#2#3#4{%
   #2\ref{#4}#3
   \catcode`\~ = 13
   \catcode`\$ = 3
   \catcode`\_ = 8
   \catcode`\# = 6
   \catcode`\& = 4
   \catcode`\% = 14
}

\@ifundefined{hyperref}{\let\hyperref\x@hyperref}{}
\@ifundefined{msihyperref}{\let\msihyperref\x@hyperref}{}




% macro for external program call
\@ifundefined{qExtProgCall}{\def\qExtProgCall#1#2#3#4#5#6{\relax}}{}
%%%%%%%%%%%%%%%%%%%%%%
%
% macros for graphics
%
\def\FILENAME#1{#1}%
%
\def\QCTOpt[#1]#2{%
  \def\QCTOptB{#1}
  \def\QCTOptA{#2}
}
\def\QCTNOpt#1{%
  \def\QCTOptA{#1}
  \let\QCTOptB\empty
}
\def\Qct{%
  \@ifnextchar[{%
    \QCTOpt}{\QCTNOpt}
}
\def\QCBOpt[#1]#2{%
  \def\QCBOptB{#1}%
  \def\QCBOptA{#2}%
}
\def\QCBNOpt#1{%
  \def\QCBOptA{#1}%
  \let\QCBOptB\empty
}
\def\Qcb{%
  \@ifnextchar[{%
    \QCBOpt}{\QCBNOpt}%
}
\def\PrepCapArgs{%
  \ifx\QCBOptA\empty
    \ifx\QCTOptA\empty
      {}%
    \else
      \ifx\QCTOptB\empty
        {\QCTOptA}%
      \else
        [\QCTOptB]{\QCTOptA}%
      \fi
    \fi
  \else
    \ifx\QCBOptA\empty
      {}%
    \else
      \ifx\QCBOptB\empty
        {\QCBOptA}%
      \else
        [\QCBOptB]{\QCBOptA}%
      \fi
    \fi
  \fi
}
\newcount\GRAPHICSTYPE
%\GRAPHICSTYPE 0 is for TurboTeX
%\GRAPHICSTYPE 1 is for DVIWindo (PostScript)
%%%(removed)%\GRAPHICSTYPE 2 is for psfig (PostScript)
\GRAPHICSTYPE=\z@
\def\GRAPHICSPS#1{%
 \ifcase\GRAPHICSTYPE%\GRAPHICSTYPE=0
   \special{ps: #1}%
 \or%\GRAPHICSTYPE=1
   \special{language "PS", include "#1"}%
%%%\or%\GRAPHICSTYPE=2
%%%  #1%
 \fi
}%
%
\def\GRAPHICSHP#1{\special{include #1}}%
%
% \graffile{ body }                                  %#1
%          { contentswidth (scalar)  }               %#2
%          { contentsheight (scalar) }               %#3
%          { vertical shift when in-line (scalar) }  %#4

\def\graffile#1#2#3#4{%
%%% \ifnum\GRAPHICSTYPE=\tw@
%%%  %Following if using psfig
%%%  \@ifundefined{psfig}{\input psfig.tex}{}%
%%%  \psfig{file=#1, height=#3, width=#2}%
%%% \else
  %Following for all others
  % JCS - added BOXTHEFRAME, see below
    \bgroup
	   \@inlabelfalse
       \leavevmode
       \@ifundefined{bbl@deactivate}{\def~{\string~}}{\activesoff}%
        \raise -#4 \BOXTHEFRAME{%
           \hbox to #2{\raise #3\hbox to #2{\null #1\hfil}}}%
    \egroup
}%
%
% A box for drafts
\def\draftbox#1#2#3#4{%
 \leavevmode\raise -#4 \hbox{%
  \frame{\rlap{\protect\tiny #1}\hbox to #2%
   {\vrule height#3 width\z@ depth\z@\hfil}%
  }%
 }%
}%
%
\newcount\@msidraft
\@msidraft=\z@
\let\nographics=\@msidraft
\newif\ifwasdraft
\wasdraftfalse

%  \GRAPHIC{ body }                                  %#1
%          { draft name }                            %#2
%          { contentswidth (scalar)  }               %#3
%          { contentsheight (scalar) }               %#4
%          { vertical shift when in-line (scalar) }  %#5
\def\GRAPHIC#1#2#3#4#5{%
   \ifnum\@msidraft=\@ne\draftbox{#2}{#3}{#4}{#5}%
   \else\graffile{#1}{#3}{#4}{#5}%
   \fi
}
%
\def\addtoLaTeXparams#1{%
    \edef\LaTeXparams{\LaTeXparams #1}}%
%
% JCS -  added a switch BoxFrame that can 
% be set by including X in the frame params.
% If set a box is drawn around the frame.

\newif\ifBoxFrame \BoxFramefalse
\newif\ifOverFrame \OverFramefalse
\newif\ifUnderFrame \UnderFramefalse

\def\BOXTHEFRAME#1{%
   \hbox{%
      \ifBoxFrame
         \frame{#1}%
      \else
         {#1}%
      \fi
   }%
}


\def\doFRAMEparams#1{\BoxFramefalse\OverFramefalse\UnderFramefalse\readFRAMEparams#1\end}%
\def\readFRAMEparams#1{%
 \ifx#1\end%
  \let\next=\relax
  \else
  \ifx#1i\dispkind=\z@\fi
  \ifx#1d\dispkind=\@ne\fi
  \ifx#1f\dispkind=\tw@\fi
  \ifx#1t\addtoLaTeXparams{t}\fi
  \ifx#1b\addtoLaTeXparams{b}\fi
  \ifx#1p\addtoLaTeXparams{p}\fi
  \ifx#1h\addtoLaTeXparams{h}\fi
  \ifx#1X\BoxFrametrue\fi
  \ifx#1O\OverFrametrue\fi
  \ifx#1U\UnderFrametrue\fi
  \ifx#1w
    \ifnum\@msidraft=1\wasdrafttrue\else\wasdraftfalse\fi
    \@msidraft=\@ne
  \fi
  \let\next=\readFRAMEparams
  \fi
 \next
 }%
%
%Macro for In-line graphics object
%   \IFRAME{ contentswidth (scalar)  }               %#1
%          { contentsheight (scalar) }               %#2
%          { vertical shift when in-line (scalar) }  %#3
%          { draft name }                            %#4
%          { body }                                  %#5
%          { caption}                                %#6


\def\IFRAME#1#2#3#4#5#6{%
      \bgroup
      \let\QCTOptA\empty
      \let\QCTOptB\empty
      \let\QCBOptA\empty
      \let\QCBOptB\empty
      #6%
      \parindent=0pt
      \leftskip=0pt
      \rightskip=0pt
      \setbox0=\hbox{\QCBOptA}%
      \@tempdima=#1\relax
      \ifOverFrame
          % Do this later
          \typeout{This is not implemented yet}%
          \show\HELP
      \else
         \ifdim\wd0>\@tempdima
            \advance\@tempdima by \@tempdima
            \ifdim\wd0 >\@tempdima
               \setbox1 =\vbox{%
                  \unskip\hbox to \@tempdima{\hfill\GRAPHIC{#5}{#4}{#1}{#2}{#3}\hfill}%
                  \unskip\hbox to \@tempdima{\parbox[b]{\@tempdima}{\QCBOptA}}%
               }%
               \wd1=\@tempdima
            \else
               \textwidth=\wd0
               \setbox1 =\vbox{%
                 \noindent\hbox to \wd0{\hfill\GRAPHIC{#5}{#4}{#1}{#2}{#3}\hfill}\\%
                 \noindent\hbox{\QCBOptA}%
               }%
               \wd1=\wd0
            \fi
         \else
            \ifdim\wd0>0pt
              \hsize=\@tempdima
              \setbox1=\vbox{%
                \unskip\GRAPHIC{#5}{#4}{#1}{#2}{0pt}%
                \break
                \unskip\hbox to \@tempdima{\hfill \QCBOptA\hfill}%
              }%
              \wd1=\@tempdima
           \else
              \hsize=\@tempdima
              \setbox1=\vbox{%
                \unskip\GRAPHIC{#5}{#4}{#1}{#2}{0pt}%
              }%
              \wd1=\@tempdima
           \fi
         \fi
         \@tempdimb=\ht1
         %\advance\@tempdimb by \dp1
         \advance\@tempdimb by -#2
         \advance\@tempdimb by #3
         \leavevmode
         \raise -\@tempdimb \hbox{\box1}%
      \fi
      \egroup%
}%
%
%Macro for Display graphics object
%   \DFRAME{ contentswidth (scalar)  }               %#1
%          { contentsheight (scalar) }               %#2
%          { draft label }                           %#3
%          { name }                                  %#4
%          { caption}                                %#5
\def\DFRAME#1#2#3#4#5{%
  \vspace\topsep
  \hfil\break
  \bgroup
     \leftskip\@flushglue
	 \rightskip\@flushglue
	 \parindent\z@
	 \parfillskip\z@skip
     \let\QCTOptA\empty
     \let\QCTOptB\empty
     \let\QCBOptA\empty
     \let\QCBOptB\empty
	 \vbox\bgroup
        \ifOverFrame 
           #5\QCTOptA\par
        \fi
        \GRAPHIC{#4}{#3}{#1}{#2}{\z@}%
        \ifUnderFrame 
           \break#5\QCBOptA
        \fi
	 \egroup
  \egroup
  \vspace\topsep
  \break
}%
%
%Macro for Floating graphic object
%   \FFRAME{ framedata f|i tbph x F|T }              %#1
%          { contentswidth (scalar)  }               %#2
%          { contentsheight (scalar) }               %#3
%          { caption }                               %#4
%          { label }                                 %#5
%          { draft name }                            %#6
%          { body }                                  %#7
\def\FFRAME#1#2#3#4#5#6#7{%
 %If float.sty loaded and float option is 'h', change to 'H'  (gp) 1998/09/05
  \@ifundefined{floatstyle}
    {%floatstyle undefined (and float.sty not present), no change
     \begin{figure}[#1]%
    }
    {%floatstyle DEFINED
	 \ifx#1h%Only the h parameter, change to H
      \begin{figure}[H]%
	 \else
      \begin{figure}[#1]%
	 \fi
	}
  \let\QCTOptA\empty
  \let\QCTOptB\empty
  \let\QCBOptA\empty
  \let\QCBOptB\empty
  \ifOverFrame
    #4
    \ifx\QCTOptA\empty
    \else
      \ifx\QCTOptB\empty
        \caption{\QCTOptA}%
      \else
        \caption[\QCTOptB]{\QCTOptA}%
      \fi
    \fi
    \ifUnderFrame\else
      \label{#5}%
    \fi
  \else
    \UnderFrametrue%
  \fi
  \begin{center}\GRAPHIC{#7}{#6}{#2}{#3}{\z@}\end{center}%
  \ifUnderFrame
    #4
    \ifx\QCBOptA\empty
      \caption{}%
    \else
      \ifx\QCBOptB\empty
        \caption{\QCBOptA}%
      \else
        \caption[\QCBOptB]{\QCBOptA}%
      \fi
    \fi
    \label{#5}%
  \fi
  \end{figure}%
 }%
%
%
%    \FRAME{ framedata f|i tbph x F|T }              %#1
%          { contentswidth (scalar)  }               %#2
%          { contentsheight (scalar) }               %#3
%          { vertical shift when in-line (scalar) }  %#4
%          { caption }                               %#5
%          { label }                                 %#6
%          { name }                                  %#7
%          { body }                                  %#8
%
%    framedata is a string which can contain the following
%    characters: idftbphxFT
%    Their meaning is as follows:
%             i, d or f : in-line, display, or floating
%             t,b,p,h   : LaTeX floating placement options
%             x         : fit contents box to contents
%             F or T    : Figure or Table. 
%                         Later this can expand
%                         to a more general float class.
%
%
\newcount\dispkind%

\def\makeactives{
  \catcode`\"=\active
  \catcode`\;=\active
  \catcode`\:=\active
  \catcode`\'=\active
  \catcode`\~=\active
}
\bgroup
   \makeactives
   \gdef\activesoff{%
      \def"{\string"}%
      \def;{\string;}%
      \def:{\string:}%
      \def'{\string'}%
      \def~{\string~}%
      %\bbl@deactivate{"}%
      %\bbl@deactivate{;}%
      %\bbl@deactivate{:}%
      %\bbl@deactivate{'}%
    }
\egroup

\def\FRAME#1#2#3#4#5#6#7#8{%
 \bgroup
 \ifnum\@msidraft=\@ne
   \wasdrafttrue
 \else
   \wasdraftfalse%
 \fi
 \def\LaTeXparams{}%
 \dispkind=\z@
 \def\LaTeXparams{}%
 \doFRAMEparams{#1}%
 \ifnum\dispkind=\z@\IFRAME{#2}{#3}{#4}{#7}{#8}{#5}\else
  \ifnum\dispkind=\@ne\DFRAME{#2}{#3}{#7}{#8}{#5}\else
   \ifnum\dispkind=\tw@
    \edef\@tempa{\noexpand\FFRAME{\LaTeXparams}}%
    \@tempa{#2}{#3}{#5}{#6}{#7}{#8}%
    \fi
   \fi
  \fi
  \ifwasdraft\@msidraft=1\else\@msidraft=0\fi{}%
  \egroup
 }%
%
% This macro added to let SW gobble a parameter that
% should not be passed on and expanded. 

\def\TEXUX#1{"texux"}

%
% Macros for text attributes:
%
\def\BF#1{{\bf {#1}}}%
\def\NEG#1{\leavevmode\hbox{\rlap{\thinspace/}{$#1$}}}%
%
%%%%%%%%%%%%%%%%%%%%%%%%%%%%%%%%%%%%%%%%%%%%%%%%%%%%%%%%%%%%%%%%%%%%%%%%
%
%
% macros for user - defined functions
\def\limfunc#1{\mathop{\rm #1}}%
\def\func#1{\mathop{\rm #1}\nolimits}%
% macro for unit names
\def\unit#1{\mathord{\thinspace\rm #1}}%

%
% miscellaneous 
\long\def\QQQ#1#2{%
     \long\expandafter\def\csname#1\endcsname{#2}}%
\@ifundefined{QTP}{\def\QTP#1{}}{}
\@ifundefined{QEXCLUDE}{\def\QEXCLUDE#1{}}{}
\@ifundefined{Qlb}{\def\Qlb#1{#1}}{}
\@ifundefined{Qlt}{\def\Qlt#1{#1}}{}
\def\QWE{}%
\long\def\QQA#1#2{}%
\def\QTR#1#2{{\csname#1\endcsname {#2}}}%
  %	Add aliases for the ulem package
  \let\QQQuline\uline
  \let\QQQsout\sout
  \let\QQQuuline\uuline
  \let\QQQuwave\uwave
  \let\QQQxout\xout
\long\def\TeXButton#1#2{#2}%
\long\def\QSubDoc#1#2{#2}%
\def\EXPAND#1[#2]#3{}%
\def\NOEXPAND#1[#2]#3{}%
\def\PROTECTED{}%
\def\LaTeXparent#1{}%
\def\ChildStyles#1{}%
\def\ChildDefaults#1{}%
\def\QTagDef#1#2#3{}%

% Constructs added with Scientific Notebook
\@ifundefined{correctchoice}{\def\correctchoice{\relax}}{}
\@ifundefined{HTML}{\def\HTML#1{\relax}}{}
\@ifundefined{TCIIcon}{\def\TCIIcon#1#2#3#4{\relax}}{}
\if@compatibility
  \typeout{Not defining UNICODE  U or CustomNote commands for LaTeX 2.09.}
\else
  \providecommand{\UNICODE}[2][]{\protect\rule{.1in}{.1in}}
  \providecommand{\U}[1]{\protect\rule{.1in}{.1in}}
  \providecommand{\CustomNote}[3][]{\marginpar{#3}}
\fi

\@ifundefined{lambdabar}{
      \def\lambdabar{\errmessage{You have used the lambdabar symbol. 
                      This is available for typesetting only in RevTeX styles.}}
   }{}

%
% Macros for style editor docs
\@ifundefined{StyleEditBeginDoc}{\def\StyleEditBeginDoc{\relax}}{}
%
% Macros for footnotes
\def\QQfnmark#1{\footnotemark}
\def\QQfntext#1#2{\addtocounter{footnote}{#1}\footnotetext{#2}}
%
% Macros for indexing.
%
\@ifundefined{TCIMAKEINDEX}{}{\makeindex}%
%
% Attempts to avoid problems with other styles
\@ifundefined{abstract}{%
 \def\abstract{%
  \if@twocolumn
   \section*{Abstract (Not appropriate in this style!)}%
   \else \small 
   \begin{center}{\bf Abstract\vspace{-.5em}\vspace{\z@}}\end{center}%
   \quotation 
   \fi
  }%
 }{%
 }%
\@ifundefined{endabstract}{\def\endabstract
  {\if@twocolumn\else\endquotation\fi}}{}%
\@ifundefined{maketitle}{\def\maketitle#1{}}{}%
\@ifundefined{affiliation}{\def\affiliation#1{}}{}%
\@ifundefined{proof}{\def\proof{\noindent{\bfseries Proof. }}}{}%
\@ifundefined{endproof}{\def\endproof{\mbox{\ \rule{.1in}{.1in}}}}{}%
\@ifundefined{newfield}{\def\newfield#1#2{}}{}%
\@ifundefined{chapter}{\def\chapter#1{\par(Chapter head:)#1\par }%
 \newcount\c@chapter}{}%
\@ifundefined{part}{\def\part#1{\par(Part head:)#1\par }}{}%
\@ifundefined{section}{\def\section#1{\par(Section head:)#1\par }}{}%
\@ifundefined{subsection}{\def\subsection#1%
 {\par(Subsection head:)#1\par }}{}%
\@ifundefined{subsubsection}{\def\subsubsection#1%
 {\par(Subsubsection head:)#1\par }}{}%
\@ifundefined{paragraph}{\def\paragraph#1%
 {\par(Subsubsubsection head:)#1\par }}{}%
\@ifundefined{subparagraph}{\def\subparagraph#1%
 {\par(Subsubsubsubsection head:)#1\par }}{}%
%%%%%%%%%%%%%%%%%%%%%%%%%%%%%%%%%%%%%%%%%%%%%%%%%%%%%%%%%%%%%%%%%%%%%%%%
% These symbols are not recognized by LaTeX
\@ifundefined{therefore}{\def\therefore{}}{}%
\@ifundefined{backepsilon}{\def\backepsilon{}}{}%
\@ifundefined{yen}{\def\yen{\hbox{\rm\rlap=Y}}}{}%
\@ifundefined{registered}{%
   \def\registered{\relax\ifmmode{}\r@gistered
                    \else$\m@th\r@gistered$\fi}%
 \def\r@gistered{^{\ooalign
  {\hfil\raise.07ex\hbox{$\scriptstyle\rm\text{R}$}\hfil\crcr
  \mathhexbox20D}}}}{}%
\@ifundefined{Eth}{\def\Eth{}}{}%
\@ifundefined{eth}{\def\eth{}}{}%
\@ifundefined{Thorn}{\def\Thorn{}}{}%
\@ifundefined{thorn}{\def\thorn{}}{}%
% A macro to allow any symbol that requires math to appear in text
\def\TEXTsymbol#1{\mbox{$#1$}}%
\@ifundefined{degree}{\def\degree{{}^{\circ}}}{}%
%
% macros for T3TeX files
\newdimen\theight
\@ifundefined{Column}{\def\Column{%
 \vadjust{\setbox\z@=\hbox{\scriptsize\quad\quad tcol}%
  \theight=\ht\z@\advance\theight by \dp\z@\advance\theight by \lineskip
  \kern -\theight \vbox to \theight{%
   \rightline{\rlap{\box\z@}}%
   \vss
   }%
  }%
 }}{}%
%
\@ifundefined{qed}{\def\qed{%
 \ifhmode\unskip\nobreak\fi\ifmmode\ifinner\else\hskip5\p@\fi\fi
 \hbox{\hskip5\p@\vrule width4\p@ height6\p@ depth1.5\p@\hskip\p@}%
 }}{}%
%
\@ifundefined{cents}{\def\cents{\hbox{\rm\rlap c/}}}{}%
\@ifundefined{tciLaplace}{\def\tciLaplace{\ensuremath{\mathcal{L}}}}{}%
\@ifundefined{tciFourier}{\def\tciFourier{\ensuremath{\mathcal{F}}}}{}%
\@ifundefined{textcurrency}{\def\textcurrency{\hbox{\rm\rlap xo}}}{}%
\@ifundefined{texteuro}{\def\texteuro{\hbox{\rm\rlap C=}}}{}%
\@ifundefined{euro}{\def\euro{\hbox{\rm\rlap C=}}}{}%
\@ifundefined{textfranc}{\def\textfranc{\hbox{\rm\rlap-F}}}{}%
\@ifundefined{textlira}{\def\textlira{\hbox{\rm\rlap L=}}}{}%
\@ifundefined{textpeseta}{\def\textpeseta{\hbox{\rm P\negthinspace s}}}{}%
%
\@ifundefined{miss}{\def\miss{\hbox{\vrule height2\p@ width 2\p@ depth\z@}}}{}%
%
\@ifundefined{vvert}{\def\vvert{\Vert}}{}%  %always translated to \left| or \right|
%
\@ifundefined{tcol}{\def\tcol#1{{\baselineskip=6\p@ \vcenter{#1}} \Column}}{}%
%
\@ifundefined{dB}{\def\dB{\hbox{{}}}}{}%        %dummy entry in column 
\@ifundefined{mB}{\def\mB#1{\hbox{$#1$}}}{}%   %column entry
\@ifundefined{nB}{\def\nB#1{\hbox{#1}}}{}%     %column entry (not math)
%
\@ifundefined{note}{\def\note{$^{\dag}}}{}%
%
\def\newfmtname{LaTeX2e}
% No longer load latexsym.  This is now handled by SWP, which uses amsfonts if necessary
%
\ifx\fmtname\newfmtname
  \DeclareOldFontCommand{\rm}{\normalfont\rmfamily}{\mathrm}
  \DeclareOldFontCommand{\sf}{\normalfont\sffamily}{\mathsf}
  \DeclareOldFontCommand{\tt}{\normalfont\ttfamily}{\mathtt}
  \DeclareOldFontCommand{\bf}{\normalfont\bfseries}{\mathbf}
  \DeclareOldFontCommand{\it}{\normalfont\itshape}{\mathit}
  \DeclareOldFontCommand{\sl}{\normalfont\slshape}{\@nomath\sl}
  \DeclareOldFontCommand{\sc}{\normalfont\scshape}{\@nomath\sc}
\fi

%
% Greek bold macros
% Redefine all of the math symbols 
% which might be bolded	 - there are 
% probably others to add to this list

\def\alpha{{\Greekmath 010B}}%
\def\beta{{\Greekmath 010C}}%
\def\gamma{{\Greekmath 010D}}%
\def\delta{{\Greekmath 010E}}%
\def\epsilon{{\Greekmath 010F}}%
\def\zeta{{\Greekmath 0110}}%
\def\eta{{\Greekmath 0111}}%
\def\theta{{\Greekmath 0112}}%
\def\iota{{\Greekmath 0113}}%
\def\kappa{{\Greekmath 0114}}%
\def\lambda{{\Greekmath 0115}}%
\def\mu{{\Greekmath 0116}}%
\def\nu{{\Greekmath 0117}}%
\def\xi{{\Greekmath 0118}}%
\def\pi{{\Greekmath 0119}}%
\def\rho{{\Greekmath 011A}}%
\def\sigma{{\Greekmath 011B}}%
\def\tau{{\Greekmath 011C}}%
\def\upsilon{{\Greekmath 011D}}%
\def\phi{{\Greekmath 011E}}%
\def\chi{{\Greekmath 011F}}%
\def\psi{{\Greekmath 0120}}%
\def\omega{{\Greekmath 0121}}%
\def\varepsilon{{\Greekmath 0122}}%
\def\vartheta{{\Greekmath 0123}}%
\def\varpi{{\Greekmath 0124}}%
\def\varrho{{\Greekmath 0125}}%
\def\varsigma{{\Greekmath 0126}}%
\def\varphi{{\Greekmath 0127}}%

\def\nabla{{\Greekmath 0272}}
\def\FindBoldGroup{%
   {\setbox0=\hbox{$\mathbf{x\global\edef\theboldgroup{\the\mathgroup}}$}}%
}

\def\Greekmath#1#2#3#4{%
    \if@compatibility
        \ifnum\mathgroup=\symbold
           \mathchoice{\mbox{\boldmath$\displaystyle\mathchar"#1#2#3#4$}}%
                      {\mbox{\boldmath$\textstyle\mathchar"#1#2#3#4$}}%
                      {\mbox{\boldmath$\scriptstyle\mathchar"#1#2#3#4$}}%
                      {\mbox{\boldmath$\scriptscriptstyle\mathchar"#1#2#3#4$}}%
        \else
           \mathchar"#1#2#3#4% 
        \fi 
    \else 
        \FindBoldGroup
        \ifnum\mathgroup=\theboldgroup % For 2e
           \mathchoice{\mbox{\boldmath$\displaystyle\mathchar"#1#2#3#4$}}%
                      {\mbox{\boldmath$\textstyle\mathchar"#1#2#3#4$}}%
                      {\mbox{\boldmath$\scriptstyle\mathchar"#1#2#3#4$}}%
                      {\mbox{\boldmath$\scriptscriptstyle\mathchar"#1#2#3#4$}}%
        \else
           \mathchar"#1#2#3#4% 
        \fi     	    
	  \fi}

\newif\ifGreekBold  \GreekBoldfalse
\let\SAVEPBF=\pbf
\def\pbf{\GreekBoldtrue\SAVEPBF}%
%

\@ifundefined{theorem}{\newtheorem{theorem}{Theorem}}{}
\@ifundefined{lemma}{\newtheorem{lemma}[theorem]{Lemma}}{}
\@ifundefined{corollary}{\newtheorem{corollary}[theorem]{Corollary}}{}
\@ifundefined{conjecture}{\newtheorem{conjecture}[theorem]{Conjecture}}{}
\@ifundefined{proposition}{\newtheorem{proposition}[theorem]{Proposition}}{}
\@ifundefined{axiom}{\newtheorem{axiom}{Axiom}}{}
\@ifundefined{remark}{\newtheorem{remark}{Remark}}{}
\@ifundefined{example}{\newtheorem{example}{Example}}{}
\@ifundefined{exercise}{\newtheorem{exercise}{Exercise}}{}
\@ifundefined{definition}{\newtheorem{definition}{Definition}}{}


\@ifundefined{mathletters}{%
  %\def\theequation{\arabic{equation}}
  \newcounter{equationnumber}  
  \def\mathletters{%
     \addtocounter{equation}{1}
     \edef\@currentlabel{\theequation}%
     \setcounter{equationnumber}{\c@equation}
     \setcounter{equation}{0}%
     \edef\theequation{\@currentlabel\noexpand\alph{equation}}%
  }
  \def\endmathletters{%
     \setcounter{equation}{\value{equationnumber}}%
  }
}{}

%Logos
\@ifundefined{BibTeX}{%
    \def\BibTeX{{\rm B\kern-.05em{\sc i\kern-.025em b}\kern-.08em
                 T\kern-.1667em\lower.7ex\hbox{E}\kern-.125emX}}}{}%
\@ifundefined{AmS}%
    {\def\AmS{{\protect\usefont{OMS}{cmsy}{m}{n}%
                A\kern-.1667em\lower.5ex\hbox{M}\kern-.125emS}}}{}%
\@ifundefined{AmSTeX}{\def\AmSTeX{\protect\AmS-\protect\TeX\@}}{}%
%

% This macro is a fix to eqnarray
\def\@@eqncr{\let\@tempa\relax
    \ifcase\@eqcnt \def\@tempa{& & &}\or \def\@tempa{& &}%
      \else \def\@tempa{&}\fi
     \@tempa
     \if@eqnsw
        \iftag@
           \@taggnum
        \else
           \@eqnnum\stepcounter{equation}%
        \fi
     \fi
     \global\tag@false
     \global\@eqnswtrue
     \global\@eqcnt\z@\cr}


\def\TCItag{\@ifnextchar*{\@TCItagstar}{\@TCItag}}
\def\@TCItag#1{%
    \global\tag@true
    \global\def\@taggnum{(#1)}%
    \global\def\@currentlabel{#1}}
\def\@TCItagstar*#1{%
    \global\tag@true
    \global\def\@taggnum{#1}%
    \global\def\@currentlabel{#1}}
%
%%%%%%%%%%%%%%%%%%%%%%%%%%%%%%%%%%%%%%%%%%%%%%%%%%%%%%%%%%%%%%%%%%%%%
%
\def\QATOP#1#2{{#1 \atop #2}}%
\def\QTATOP#1#2{{\textstyle {#1 \atop #2}}}%
\def\QDATOP#1#2{{\displaystyle {#1 \atop #2}}}%
\def\QABOVE#1#2#3{{#2 \above#1 #3}}%
\def\QTABOVE#1#2#3{{\textstyle {#2 \above#1 #3}}}%
\def\QDABOVE#1#2#3{{\displaystyle {#2 \above#1 #3}}}%
\def\QOVERD#1#2#3#4{{#3 \overwithdelims#1#2 #4}}%
\def\QTOVERD#1#2#3#4{{\textstyle {#3 \overwithdelims#1#2 #4}}}%
\def\QDOVERD#1#2#3#4{{\displaystyle {#3 \overwithdelims#1#2 #4}}}%
\def\QATOPD#1#2#3#4{{#3 \atopwithdelims#1#2 #4}}%
\def\QTATOPD#1#2#3#4{{\textstyle {#3 \atopwithdelims#1#2 #4}}}%
\def\QDATOPD#1#2#3#4{{\displaystyle {#3 \atopwithdelims#1#2 #4}}}%
\def\QABOVED#1#2#3#4#5{{#4 \abovewithdelims#1#2#3 #5}}%
\def\QTABOVED#1#2#3#4#5{{\textstyle 
   {#4 \abovewithdelims#1#2#3 #5}}}%
\def\QDABOVED#1#2#3#4#5{{\displaystyle 
   {#4 \abovewithdelims#1#2#3 #5}}}%
%
% Macros for text size operators:
%

\def\tint{\msi@int\textstyle\int}%
\def\tiint{\msi@int\textstyle\iint}%
\def\tiiint{\msi@int\textstyle\iiint}%
\def\tiiiint{\msi@int\textstyle\iiiint}%
\def\tidotsint{\msi@int\textstyle\idotsint}%
\def\toint{\msi@int\textstyle\oint}%


\def\tsum{\mathop{\textstyle \sum }}%
\def\tprod{\mathop{\textstyle \prod }}%
\def\tbigcap{\mathop{\textstyle \bigcap }}%
\def\tbigwedge{\mathop{\textstyle \bigwedge }}%
\def\tbigoplus{\mathop{\textstyle \bigoplus }}%
\def\tbigodot{\mathop{\textstyle \bigodot }}%
\def\tbigsqcup{\mathop{\textstyle \bigsqcup }}%
\def\tcoprod{\mathop{\textstyle \coprod }}%
\def\tbigcup{\mathop{\textstyle \bigcup }}%
\def\tbigvee{\mathop{\textstyle \bigvee }}%
\def\tbigotimes{\mathop{\textstyle \bigotimes }}%
\def\tbiguplus{\mathop{\textstyle \biguplus }}%
%
%
%Macros for display size operators:
%

\newtoks\temptoksa
\newtoks\temptoksb
\newtoks\temptoksc


\def\msi@int#1#2{%
 \def\@temp{{#1#2\the\temptoksc_{\the\temptoksa}^{\the\temptoksb}}}%   
 \futurelet\@nextcs
 \@int
}

\def\@int{%
   \ifx\@nextcs\limits
      \typeout{Found limits}%
      \temptoksc={\limits}%
	  \let\@next\@intgobble%
   \else\ifx\@nextcs\nolimits
      \typeout{Found nolimits}%
      \temptoksc={\nolimits}%
	  \let\@next\@intgobble%
   \else
      \typeout{Did not find limits or no limits}%
      \temptoksc={}%
      \let\@next\msi@limits%
   \fi\fi
   \@next   
}%

\def\@intgobble#1{%
   \typeout{arg is #1}%
   \msi@limits
}


\def\msi@limits{%
   \temptoksa={}%
   \temptoksb={}%
   \@ifnextchar_{\@limitsa}{\@limitsb}%
}

\def\@limitsa_#1{%
   \temptoksa={#1}%
   \@ifnextchar^{\@limitsc}{\@temp}%
}

\def\@limitsb{%
   \@ifnextchar^{\@limitsc}{\@temp}%
}

\def\@limitsc^#1{%
   \temptoksb={#1}%
   \@ifnextchar_{\@limitsd}{\@temp}%   
}

\def\@limitsd_#1{%
   \temptoksa={#1}%
   \@temp
}



\def\dint{\msi@int\displaystyle\int}%
\def\diint{\msi@int\displaystyle\iint}%
\def\diiint{\msi@int\displaystyle\iiint}%
\def\diiiint{\msi@int\displaystyle\iiiint}%
\def\didotsint{\msi@int\displaystyle\idotsint}%
\def\doint{\msi@int\displaystyle\oint}%

\def\dsum{\mathop{\displaystyle \sum }}%
\def\dprod{\mathop{\displaystyle \prod }}%
\def\dbigcap{\mathop{\displaystyle \bigcap }}%
\def\dbigwedge{\mathop{\displaystyle \bigwedge }}%
\def\dbigoplus{\mathop{\displaystyle \bigoplus }}%
\def\dbigodot{\mathop{\displaystyle \bigodot }}%
\def\dbigsqcup{\mathop{\displaystyle \bigsqcup }}%
\def\dcoprod{\mathop{\displaystyle \coprod }}%
\def\dbigcup{\mathop{\displaystyle \bigcup }}%
\def\dbigvee{\mathop{\displaystyle \bigvee }}%
\def\dbigotimes{\mathop{\displaystyle \bigotimes }}%
\def\dbiguplus{\mathop{\displaystyle \biguplus }}%

\if@compatibility\else
  % Always load amsmath in LaTeX2e mode
  \RequirePackage{amsmath}
\fi

\def\ExitTCILatex{\makeatother\endinput}

\bgroup
\ifx\ds@amstex\relax
   \message{amstex already loaded}\aftergroup\ExitTCILatex
\else
   \@ifpackageloaded{amsmath}%
      {\if@compatibility\message{amsmath already loaded}\fi\aftergroup\ExitTCILatex}
      {}
   \@ifpackageloaded{amstex}%
      {\if@compatibility\message{amstex already loaded}\fi\aftergroup\ExitTCILatex}
      {}
   \@ifpackageloaded{amsgen}%
      {\if@compatibility\message{amsgen already loaded}\fi\aftergroup\ExitTCILatex}
      {}
\fi
\egroup

%Exit if any of the AMS macros are already loaded.
%This is always the case for LaTeX2e mode.


%%%%%%%%%%%%%%%%%%%%%%%%%%%%%%%%%%%%%%%%%%%%%%%%%%%%%%%%%%%%%%%%%%%%%%%%%%
% NOTE: The rest of this file is read only if in LaTeX 2.09 compatibility
% mode. This section is used to define AMS-like constructs in the
% event they have not been defined.
%%%%%%%%%%%%%%%%%%%%%%%%%%%%%%%%%%%%%%%%%%%%%%%%%%%%%%%%%%%%%%%%%%%%%%%%%%
\typeout{TCILATEX defining AMS-like constructs in LaTeX 2.09 COMPATIBILITY MODE}
%%%%%%%%%%%%%%%%%%%%%%%%%%%%%%%%%%%%%%%%%%%%%%%%%%%%%%%%%%%%%%%%%%%%%%%%
%  Macros to define some AMS LaTeX constructs when 
%  AMS LaTeX has not been loaded
% 
% These macros are copied from the AMS-TeX package for doing
% multiple integrals.
%
\let\DOTSI\relax
\def\RIfM@{\relax\ifmmode}%
\def\FN@{\futurelet\next}%
\newcount\intno@
\def\iint{\DOTSI\intno@\tw@\FN@\ints@}%
\def\iiint{\DOTSI\intno@\thr@@\FN@\ints@}%
\def\iiiint{\DOTSI\intno@4 \FN@\ints@}%
\def\idotsint{\DOTSI\intno@\z@\FN@\ints@}%
\def\ints@{\findlimits@\ints@@}%
\newif\iflimtoken@
\newif\iflimits@
\def\findlimits@{\limtoken@true\ifx\next\limits\limits@true
 \else\ifx\next\nolimits\limits@false\else
 \limtoken@false\ifx\ilimits@\nolimits\limits@false\else
 \ifinner\limits@false\else\limits@true\fi\fi\fi\fi}%
\def\multint@{\int\ifnum\intno@=\z@\intdots@                          %1
 \else\intkern@\fi                                                    %2
 \ifnum\intno@>\tw@\int\intkern@\fi                                   %3
 \ifnum\intno@>\thr@@\int\intkern@\fi                                 %4
 \int}%                                                               %5
\def\multintlimits@{\intop\ifnum\intno@=\z@\intdots@\else\intkern@\fi
 \ifnum\intno@>\tw@\intop\intkern@\fi
 \ifnum\intno@>\thr@@\intop\intkern@\fi\intop}%
\def\intic@{%
    \mathchoice{\hskip.5em}{\hskip.4em}{\hskip.4em}{\hskip.4em}}%
\def\negintic@{\mathchoice
 {\hskip-.5em}{\hskip-.4em}{\hskip-.4em}{\hskip-.4em}}%
\def\ints@@{\iflimtoken@                                              %1
 \def\ints@@@{\iflimits@\negintic@
   \mathop{\intic@\multintlimits@}\limits                             %2
  \else\multint@\nolimits\fi                                          %3
  \eat@}%                                                             %4
 \else                                                                %5
 \def\ints@@@{\iflimits@\negintic@
  \mathop{\intic@\multintlimits@}\limits\else
  \multint@\nolimits\fi}\fi\ints@@@}%
\def\intkern@{\mathchoice{\!\!\!}{\!\!}{\!\!}{\!\!}}%
\def\plaincdots@{\mathinner{\cdotp\cdotp\cdotp}}%
\def\intdots@{\mathchoice{\plaincdots@}%
 {{\cdotp}\mkern1.5mu{\cdotp}\mkern1.5mu{\cdotp}}%
 {{\cdotp}\mkern1mu{\cdotp}\mkern1mu{\cdotp}}%
 {{\cdotp}\mkern1mu{\cdotp}\mkern1mu{\cdotp}}}%
%
%
%  These macros are for doing the AMS \text{} construct
%
\def\RIfM@{\relax\protect\ifmmode}
\def\text{\RIfM@\expandafter\text@\else\expandafter\mbox\fi}
\let\nfss@text\text
\def\text@#1{\mathchoice
   {\textdef@\displaystyle\f@size{#1}}%
   {\textdef@\textstyle\tf@size{\firstchoice@false #1}}%
   {\textdef@\textstyle\sf@size{\firstchoice@false #1}}%
   {\textdef@\textstyle \ssf@size{\firstchoice@false #1}}%
   \glb@settings}

\def\textdef@#1#2#3{\hbox{{%
                    \everymath{#1}%
                    \let\f@size#2\selectfont
                    #3}}}
\newif\iffirstchoice@
\firstchoice@true
%
%These are the AMS constructs for multiline limits.
%
\def\Let@{\relax\iffalse{\fi\let\\=\cr\iffalse}\fi}%
\def\vspace@{\def\vspace##1{\crcr\noalign{\vskip##1\relax}}}%
\def\multilimits@{\bgroup\vspace@\Let@
 \baselineskip\fontdimen10 \scriptfont\tw@
 \advance\baselineskip\fontdimen12 \scriptfont\tw@
 \lineskip\thr@@\fontdimen8 \scriptfont\thr@@
 \lineskiplimit\lineskip
 \vbox\bgroup\ialign\bgroup\hfil$\m@th\scriptstyle{##}$\hfil\crcr}%
\def\Sb{_\multilimits@}%
\def\endSb{\crcr\egroup\egroup\egroup}%
\def\Sp{^\multilimits@}%
\let\endSp\endSb
%
%
%These are AMS constructs for horizontal arrows
%
\newdimen\ex@
\ex@.2326ex
\def\rightarrowfill@#1{$#1\m@th\mathord-\mkern-6mu\cleaders
 \hbox{$#1\mkern-2mu\mathord-\mkern-2mu$}\hfill
 \mkern-6mu\mathord\rightarrow$}%
\def\leftarrowfill@#1{$#1\m@th\mathord\leftarrow\mkern-6mu\cleaders
 \hbox{$#1\mkern-2mu\mathord-\mkern-2mu$}\hfill\mkern-6mu\mathord-$}%
\def\leftrightarrowfill@#1{$#1\m@th\mathord\leftarrow
\mkern-6mu\cleaders
 \hbox{$#1\mkern-2mu\mathord-\mkern-2mu$}\hfill
 \mkern-6mu\mathord\rightarrow$}%
\def\overrightarrow{\mathpalette\overrightarrow@}%
\def\overrightarrow@#1#2{\vbox{\ialign{##\crcr\rightarrowfill@#1\crcr
 \noalign{\kern-\ex@\nointerlineskip}$\m@th\hfil#1#2\hfil$\crcr}}}%
\let\overarrow\overrightarrow
\def\overleftarrow{\mathpalette\overleftarrow@}%
\def\overleftarrow@#1#2{\vbox{\ialign{##\crcr\leftarrowfill@#1\crcr
 \noalign{\kern-\ex@\nointerlineskip}$\m@th\hfil#1#2\hfil$\crcr}}}%
\def\overleftrightarrow{\mathpalette\overleftrightarrow@}%
\def\overleftrightarrow@#1#2{\vbox{\ialign{##\crcr
   \leftrightarrowfill@#1\crcr
 \noalign{\kern-\ex@\nointerlineskip}$\m@th\hfil#1#2\hfil$\crcr}}}%
\def\underrightarrow{\mathpalette\underrightarrow@}%
\def\underrightarrow@#1#2{\vtop{\ialign{##\crcr$\m@th\hfil#1#2\hfil
  $\crcr\noalign{\nointerlineskip}\rightarrowfill@#1\crcr}}}%
\let\underarrow\underrightarrow
\def\underleftarrow{\mathpalette\underleftarrow@}%
\def\underleftarrow@#1#2{\vtop{\ialign{##\crcr$\m@th\hfil#1#2\hfil
  $\crcr\noalign{\nointerlineskip}\leftarrowfill@#1\crcr}}}%
\def\underleftrightarrow{\mathpalette\underleftrightarrow@}%
\def\underleftrightarrow@#1#2{\vtop{\ialign{##\crcr$\m@th
  \hfil#1#2\hfil$\crcr
 \noalign{\nointerlineskip}\leftrightarrowfill@#1\crcr}}}%
%%%%%%%%%%%%%%%%%%%%%

\def\qopnamewl@#1{\mathop{\operator@font#1}\nlimits@}
\let\nlimits@\displaylimits
\def\setboxz@h{\setbox\z@\hbox}


\def\varlim@#1#2{\mathop{\vtop{\ialign{##\crcr
 \hfil$#1\m@th\operator@font lim$\hfil\crcr
 \noalign{\nointerlineskip}#2#1\crcr
 \noalign{\nointerlineskip\kern-\ex@}\crcr}}}}

 \def\rightarrowfill@#1{\m@th\setboxz@h{$#1-$}\ht\z@\z@
  $#1\copy\z@\mkern-6mu\cleaders
  \hbox{$#1\mkern-2mu\box\z@\mkern-2mu$}\hfill
  \mkern-6mu\mathord\rightarrow$}
\def\leftarrowfill@#1{\m@th\setboxz@h{$#1-$}\ht\z@\z@
  $#1\mathord\leftarrow\mkern-6mu\cleaders
  \hbox{$#1\mkern-2mu\copy\z@\mkern-2mu$}\hfill
  \mkern-6mu\box\z@$}


\def\projlim{\qopnamewl@{proj\,lim}}
\def\injlim{\qopnamewl@{inj\,lim}}
\def\varinjlim{\mathpalette\varlim@\rightarrowfill@}
\def\varprojlim{\mathpalette\varlim@\leftarrowfill@}
\def\varliminf{\mathpalette\varliminf@{}}
\def\varliminf@#1{\mathop{\underline{\vrule\@depth.2\ex@\@width\z@
   \hbox{$#1\m@th\operator@font lim$}}}}
\def\varlimsup{\mathpalette\varlimsup@{}}
\def\varlimsup@#1{\mathop{\overline
  {\hbox{$#1\m@th\operator@font lim$}}}}

%
%Companion to stackrel
\def\stackunder#1#2{\mathrel{\mathop{#2}\limits_{#1}}}%
%
%
% These are AMS environments that will be defined to
% be verbatims if amstex has not actually been 
% loaded
%
%
\begingroup \catcode `|=0 \catcode `[= 1
\catcode`]=2 \catcode `\{=12 \catcode `\}=12
\catcode`\\=12 
|gdef|@alignverbatim#1\end{align}[#1|end[align]]
|gdef|@salignverbatim#1\end{align*}[#1|end[align*]]

|gdef|@alignatverbatim#1\end{alignat}[#1|end[alignat]]
|gdef|@salignatverbatim#1\end{alignat*}[#1|end[alignat*]]

|gdef|@xalignatverbatim#1\end{xalignat}[#1|end[xalignat]]
|gdef|@sxalignatverbatim#1\end{xalignat*}[#1|end[xalignat*]]

|gdef|@gatherverbatim#1\end{gather}[#1|end[gather]]
|gdef|@sgatherverbatim#1\end{gather*}[#1|end[gather*]]

|gdef|@gatherverbatim#1\end{gather}[#1|end[gather]]
|gdef|@sgatherverbatim#1\end{gather*}[#1|end[gather*]]


|gdef|@multilineverbatim#1\end{multiline}[#1|end[multiline]]
|gdef|@smultilineverbatim#1\end{multiline*}[#1|end[multiline*]]

|gdef|@arraxverbatim#1\end{arrax}[#1|end[arrax]]
|gdef|@sarraxverbatim#1\end{arrax*}[#1|end[arrax*]]

|gdef|@tabulaxverbatim#1\end{tabulax}[#1|end[tabulax]]
|gdef|@stabulaxverbatim#1\end{tabulax*}[#1|end[tabulax*]]


|endgroup
  

  
\def\align{\@verbatim \frenchspacing\@vobeyspaces \@alignverbatim
You are using the "align" environment in a style in which it is not defined.}
\let\endalign=\endtrivlist
 
\@namedef{align*}{\@verbatim\@salignverbatim
You are using the "align*" environment in a style in which it is not defined.}
\expandafter\let\csname endalign*\endcsname =\endtrivlist




\def\alignat{\@verbatim \frenchspacing\@vobeyspaces \@alignatverbatim
You are using the "alignat" environment in a style in which it is not defined.}
\let\endalignat=\endtrivlist
 
\@namedef{alignat*}{\@verbatim\@salignatverbatim
You are using the "alignat*" environment in a style in which it is not defined.}
\expandafter\let\csname endalignat*\endcsname =\endtrivlist




\def\xalignat{\@verbatim \frenchspacing\@vobeyspaces \@xalignatverbatim
You are using the "xalignat" environment in a style in which it is not defined.}
\let\endxalignat=\endtrivlist
 
\@namedef{xalignat*}{\@verbatim\@sxalignatverbatim
You are using the "xalignat*" environment in a style in which it is not defined.}
\expandafter\let\csname endxalignat*\endcsname =\endtrivlist




\def\gather{\@verbatim \frenchspacing\@vobeyspaces \@gatherverbatim
You are using the "gather" environment in a style in which it is not defined.}
\let\endgather=\endtrivlist
 
\@namedef{gather*}{\@verbatim\@sgatherverbatim
You are using the "gather*" environment in a style in which it is not defined.}
\expandafter\let\csname endgather*\endcsname =\endtrivlist


\def\multiline{\@verbatim \frenchspacing\@vobeyspaces \@multilineverbatim
You are using the "multiline" environment in a style in which it is not defined.}
\let\endmultiline=\endtrivlist
 
\@namedef{multiline*}{\@verbatim\@smultilineverbatim
You are using the "multiline*" environment in a style in which it is not defined.}
\expandafter\let\csname endmultiline*\endcsname =\endtrivlist


\def\arrax{\@verbatim \frenchspacing\@vobeyspaces \@arraxverbatim
You are using a type of "array" construct that is only allowed in AmS-LaTeX.}
\let\endarrax=\endtrivlist

\def\tabulax{\@verbatim \frenchspacing\@vobeyspaces \@tabulaxverbatim
You are using a type of "tabular" construct that is only allowed in AmS-LaTeX.}
\let\endtabulax=\endtrivlist

 
\@namedef{arrax*}{\@verbatim\@sarraxverbatim
You are using a type of "array*" construct that is only allowed in AmS-LaTeX.}
\expandafter\let\csname endarrax*\endcsname =\endtrivlist

\@namedef{tabulax*}{\@verbatim\@stabulaxverbatim
You are using a type of "tabular*" construct that is only allowed in AmS-LaTeX.}
\expandafter\let\csname endtabulax*\endcsname =\endtrivlist

% macro to simulate ams tag construct


% This macro is a fix to the equation environment
 \def\endequation{%
     \ifmmode\ifinner % FLEQN hack
      \iftag@
        \addtocounter{equation}{-1} % undo the increment made in the begin part
        $\hfil
           \displaywidth\linewidth\@taggnum\egroup \endtrivlist
        \global\tag@false
        \global\@ignoretrue   
      \else
        $\hfil
           \displaywidth\linewidth\@eqnnum\egroup \endtrivlist
        \global\tag@false
        \global\@ignoretrue 
      \fi
     \else   
      \iftag@
        \addtocounter{equation}{-1} % undo the increment made in the begin part
        \eqno \hbox{\@taggnum}
        \global\tag@false%
        $$\global\@ignoretrue
      \else
        \eqno \hbox{\@eqnnum}% $$ BRACE MATCHING HACK
        $$\global\@ignoretrue
      \fi
     \fi\fi
 } 

 \newif\iftag@ \tag@false
 
 \def\TCItag{\@ifnextchar*{\@TCItagstar}{\@TCItag}}
 \def\@TCItag#1{%
     \global\tag@true
     \global\def\@taggnum{(#1)}%
     \global\def\@currentlabel{#1}}
 \def\@TCItagstar*#1{%
     \global\tag@true
     \global\def\@taggnum{#1}%
     \global\def\@currentlabel{#1}}

  \@ifundefined{tag}{
     \def\tag{\@ifnextchar*{\@tagstar}{\@tag}}
     \def\@tag#1{%
         \global\tag@true
         \global\def\@taggnum{(#1)}}
     \def\@tagstar*#1{%
         \global\tag@true
         \global\def\@taggnum{#1}}
  }{}

\def\tfrac#1#2{{\textstyle {#1 \over #2}}}%
\def\dfrac#1#2{{\displaystyle {#1 \over #2}}}%
\def\binom#1#2{{#1 \choose #2}}%
\def\tbinom#1#2{{\textstyle {#1 \choose #2}}}%
\def\dbinom#1#2{{\displaystyle {#1 \choose #2}}}%

% Do not add anything to the end of this file.  
% The last section of the file is loaded only if 
% amstex has not been.
\makeatother
\endinput

\begin{document}

\title{\textsc{Marginal Effects for Probit and Tobit with Endogeneity}}
\author{Kirill S. \textsc{Evdokimov}\thanks{%
Universitat Pompeu Fabra and Barcelona School of Economics: \textsf{%
kirill.evdokimov@upf.edu}.} \and Ilze \textsc{Kalnina}\thanks{%
North Carolina State University: \textsf{ikalnin@ncsu.edu}.} \and Andrei 
\textsc{Zeleneev}\thanks{%
University College London: \textsf{a.zeleneev@ucl.ac.uk}.} \thanks{
Evdokimov gratefully acknowledges the support from the Spanish MCIN/AEI via
grants RYC2020-030623-I, PID2019-107352GB-I00, and Severo Ochoa Programme
CEX2019-000915-S.}}
\date{This version: June 22, 2023} %
\maketitle

\begin{abstract}

When evaluating partial effects, it is important to distinguish between
structural endogeneity and measurement errors. In contrast to linear models,
these two sources of endogeneity affect partial effects differently in
nonlinear models. We study this issue focusing on the Instrumental Variable
(IV) Probit and Tobit models. We show that even when a valid IV is
available, failing to differentiate between the two types of endogeneity can
lead to either under- or over-estimation of the partial effects. We develop
simple estimators of the bounds on the partial effects and provide easy to
implement confidence intervals that correctly account for both types of
endogeneity. We illustrate the methods in a Monte Carlo simulation and an
empirical application.

\bigskip

Keywords: (Average) Partial Effects, Instrumental Variable, Control
Variable, Errors-in-Variables, Counterfactuals
\end{abstract}

\newpage

\section{Introduction}

Probit and Tobit are some of the most popular nonlinear models in applied
economics. When a covariate is endogenous, IV-Probit and IV-Tobit models are
used for instrumental variable (IV) estimation of the coefficients (%
\citealp{SmithBlundell1986Ecta}, \citealp{RiversVuong1988JoE}).\footnote{%
For example, in Stata, these estimators are \textit{ivprobit} and \textit{%
ivtobit}.}

A covariate can be endogenous for two reasons. First, the covariate can be
correlated with the individual's unobserved characteristics (unobserved
heterogeneity). Second, mismeasurement of the covariate
(Errors-in-Variables, EiV) also results in endogeneity. We will refer to
these two types of endogeneity as the \textquotedblleft
structural\textquotedblright\ endogeneity and the EiV. In many empirical
settings both sources of endogeneity need to be addressed simultaneously.

The goal of this paper is to characterize the partial effects in the classic
IV-Probit and IV-Tobit models, allowing for both types of endogeneity, and
to emphasize the importance of distinguishing between the two types. We
provide the expressions for the partial effects and average partial effects
that correctly account for the two kinds of endogeneity. Although the two
sources of endogeneity cannot be precisely distinguished using the observed
data, we use the constraints of the model to obtain bounds on the amounts of
endogeneity that can be attributed to each source. We use these bounds to
provide simple estimators of the bounds on the partial and average partial
effects, allowing for both types of endogeneity. We also provide
corresponding valid confidence intervals that are easy to calculate.

In nonlinear models, the need to differentiate between the two kinds of
endogeneity arises because structural endogeneity and EiV play different
roles. In particular, partial effects of covariates are averaged with
respect to the distribution of the individual unobserved heterogeneity. On
the other hand, one aims to remove the impact of the measurement errors,
since they are not properties of individuals but a deficiency in the
measurement process. Thus, even though the IV-Probit and IV-Tobit methods
consistently estimate the coefficients on all regressors regardless of the
source of endogeneity, the effects of the covariates on the outcomes are
only partially identified, since the distribution of the unobserved
heterogeneity is partially identified. The width of the identified set
depends on how hard it is to disentangle structural endogeneity and EiV for
the data at hand. Importantly, we find that naively ignoring the distinction
between the two types of endogeneity can result in both under- and
over-estimation of the magnitude of the partial effects by these IV
estimators.\footnote{\cite{Wooldridge2010Book}, page 586, alludes to the
potential importance of the sources of endogeneity for the partial effects
in IV-Probit, but does not elaborate.}

IV-Probit and IV-Tobit can be interpreted as control variable estimators.
Partial effects in general control variable models were considered by \cite%
{BlundellPowell2003Advances}, \cite{Chesher2003Ecta}, \cite%
{ImbensNewey2009Ecta}, and \cite{Wooldridge2005APE,Wooldridge2015JHR}, among
others. In nonlinear models, accounting for both types of endogeneity is
difficult, see, e.g., \cite{Schennach2022JEcLit}, and it is not clear how to
characterize the bounds on the partial effects in practice. Recently, \cite%
{ChesherKimRosen2023JoE} studied IV methods for the Tobit models relaxing
the first stage specification assumptions, but did not consider EiV.

The assumption of gaussianity of the unobservables simplifies the analysis
of the IV-Probit and IV-Tobit models and allows us to obtain simple bounds
and confidence intervals on the true partial effects that are easy to
implement. It would be important to relax the gaussianity assumption in the
future. At the same time, the simplicity of the proposed approach makes it a
convenient starting point in empirical analysis. 
In particular, it allows researchers to gauge the importance of properly
accounting for both types of endogeneity, which is essential given the
ubiquity of both in economics.

The rest of the paper is organized as follows. The analysis of partial
effects in the Probit and Tobit models is virtually identical, thus we first
consider Tobit in Sections \ref{sec:Model}-\ref{sec:Analysis}. In Section~%
\ref{sec:Extensions}, we explain how the methods apply to the Probit model,
and extend the analysis to cover the average partial effects and other
counterfactuals. Section~\ref{sec:MC} provides some Monte Carlo simulation
results. Section~\ref{sec:Empirical} presents an illustrative empirical
application.

\section{The Model\label{sec:Model}}

The Tobit model is often used for estimation of economic models with a
\textquotedblleft corner solution,\textquotedblright\ i.e., models where the
outcome variable $Y_{i}\,$\ is forced to be non-negative. The examples of
such dependent variables $Y_{i}$ include the amounts of charitable
contributions, hours worked, or monthly consumption of cigarettes.

First, consider the standard Tobit model with exogenous covariates and
without EiV:%
\begin{equation}
Y_{i}=m\left( \theta _{01}X_{i}^{\ast }+\theta _{02}^{\prime
}W_{i}+U_{i}^{\ast }\right) ,\text{\quad where }m\left( s\right) =\max (s,0),
\label{eq:dgp Y exog copy}
\end{equation}%
the individual unobserved heterogeneity $U_{i}^{\ast }$ has a normal
distribution $N\left( 0,\sigma _{U^{\ast }}^{2}\right) $ and is independent
from the covariates $X_{i}^{\ast }$ and $W_{i}$. We use the asterisk to
denote variables that will be affected by the EiV, as we explain in detail
below.{}

We collect the covariates in a vector $H_{i}^{\ast }=\left( X_{i}^{\ast
},W_{i}^{\prime }\right) ^{\prime }$, so (\ref{eq:dgp Y exog copy}) can be
written as $Y_{i}=m\left( \theta _{0}^{\prime }H_{i}^{\ast }+U_{i}^{\ast
}\right) $. The standard normal cumulative distribution and density
functions are denoted by $\Phi $ and $\phi $, respectively.\footnote{%
Most of the analysis in Sections~\ref{sec:Model} and \ref{sec:Analysis}
equally applies to the Probit model. For simplicity of exposition, we focus
on the Tobit model for the moment, and then discuss Probit in Section~\ref%
{sec:Probit}.}

In the Tobit model, one is usually interested in the partial effects
(marginal effects) of covariates $H_{i}^{\ast }$ on $\mathrm{E}\left(
Y_{i}|H_{i}^{\ast }\right) $ and $\mathrm{P}\left( Y_{i}>0|H_{i}^{\ast
}\right) $. For concreteness we consider partial effects of the continuously
distributed covariates.

The partial effect of the $j^{th}$ covariate on the mean {}$\mathrm{E}\left( Y_{i}|H_{i}^{\ast }=h\right) $ at a
given $h$ is%
\begin{equation}
PE_{j}^{\text{Tob}}\left( h\right) =\frac{\partial }{\partial h_{j}}\mathrm{E%
}\left( Y_{i}|H_{i}^{\ast }=h\right) =\Phi \left( \frac{\theta _{0}^{\prime
}h}{\sigma _{U^{\ast }}}\right) \theta _{0j}.  \label{eq:PE_Tobit}
\end{equation}%
The partial effect of the $j^{th}$ covariate on the probability $P\left(
Y_{i}>0|H_{i}^{\ast }=h\right) $ is 
\begin{equation}
PE_{j}^{\Pr }\left( h\right) =\frac{\partial }{\partial h_{j}}\mathrm{P}%
\left( Y_{i}>0|H_{i}^{\ast }=h\right) =\phi \left( \frac{\theta _{0}^{\prime
}h}{\sigma _{U^{\ast }}}\right) \frac{\theta _{0j}}{\sigma _{U^{\ast }}}.
\label{eq:PE_Probit}
\end{equation}%
These formulas for the $PE_{j}$ are standard, see, e.g., \cite%
{Wooldridge2010Book}, for detailed calculations. Most often one considers
the partial effects at the means of the covariates $h=E\left[ H_{i}^{\ast }%
\right] $.

When $X_{i}^{\ast }$ is correlated with $U_{i}^{\ast }$ and we observe data $%
\left( Y_{i},X_{i}^{\ast },W_{i},Z_{i}\right) $, the IV-Tobit model can be
estimated using instrumental variables $Z_{i}$, as proposed by \cite%
{SmithBlundell1986Ecta}, \cite{Newey1987}, and \cite{RiversVuong1988JoE}.
Assume that%
\begin{eqnarray}
Y_{i} &=&m\left( \theta _{01}X_{i}^{\ast }+\theta _{02}^{\prime
}W_{i}+U_{i}^{\ast }\right) ,\text{\quad }m\left( s\right) =\max (s,0),
\label{eq:dgp Y exog} \\
X_{i}^{\ast } &=&\pi _{01}^{\prime }Z_{i}+\pi _{02}^{\prime
}W_{i}+V_{i}^{\ast },  \label{eq:dgp Xstar gen}
\end{eqnarray}%
where $V_{i}^{\ast }$ is a normal random variable, possibly correlated with $%
U_{i}^{\ast }$:%
\begin{equation}
\left( 
\begin{array}{c}
U_{i}^{\ast } \\ 
V_{i}^{\ast }%
\end{array}%
\right) \sim N\left( \left( 
\begin{array}{c}
0 \\ 
0%
\end{array}%
\right) ,%
\begin{pmatrix}
\sigma _{U^{\ast }}^{2} & \sigma _{U^{\ast }V^{\ast }} \\ 
\sigma _{U^{\ast }V^{\ast }} & \sigma _{V^{\ast }}^{2}%
\end{pmatrix}%
\right) .  \label{eq:dgp UV gen}
\end{equation}

The IV-Tobit model in (\ref{eq:dgp Y exog})-(\ref{eq:dgp UV gen}) can be
estimated using a random sample of $\left( Y_{i},X_{i}^{\ast
},W_{i},Z_{i}\right) $ in two steps, see, e.g., \cite{Wooldridge2010Book}.
First, one estimates $V_{i}^{\ast }$ in equation (\ref{eq:dgp Xstar gen}) by
the the residuals $\hat{V}_{i}^{\ast }$ in the regression of $X_{i}^{\ast }$
on $\left( W_{i},Z_{i}\right) $. Note that we can write $U_{i}^{\ast
}=e_{i}^{\ast }+\theta _{V^{\ast }}V_{i}^{\ast }$, where $\theta _{V^{\ast
}}\equiv \sigma _{V^{\ast }U^{\ast }}/\sigma _{V^{\ast }}^{2}$, and $%
e_{i}^{\ast }$ is independent of $Z_{i}$, $W_{i}$, and $V_{i}^{\ast }$ (and
hence of $X_{i}^{\ast }$). Then, one estimates the standard Tobit model 
\begin{equation*}
Y_{i}=m\left( \theta _{01}X_{i}^{\ast }+\theta _{02}^{\prime }W_{i}+\theta
_{V^{\ast }}V_{i}^{\ast }+e_{i}^{\ast }\right) ,
\end{equation*}%
where $V_{i}^{\ast }$ are replaced by their estimates $\hat{V}_{i}^{\ast }$.
(Alternatively, the two steps can be combined and all of the parameters can
be estimated simultaneously by the Maximum Likelihood Estimator.) The reason
this approach works is that equation (\ref{eq:dgp Xstar gen}) creates a
control variable $V_{i}^{\ast }$, and the inclusion of $V_{i}^{\ast }$ in
the above equation makes $X_{i}^{\ast }$ exogenous. {}

To estimate the partial effects, one would plug the estimates $\widehat{%
\theta }$ and $\widehat{\sigma }_{U^{\ast }}^{2}$ into equations~(\ref%
{eq:PE_Tobit})-(\ref{eq:PE_Probit}) in place of $\theta _{0}$ and $\sigma
_{U^{\ast }}^{2}$.

So far we were assuming that the data has no measurement errors. We now
allow $X_{i}^{\ast }$ to be mismeasured, i.e., that instead of $X_{i}^{\ast
} $ we observe its noisy measurement $X_{i}$:%
\begin{equation}
X_{i}=X_{i}^{\ast }+\varepsilon _{i},\qquad \varepsilon _{i}\sim N\left(
0,\sigma _{\varepsilon }^{2}\right) .  \label{eq:ME}
\end{equation}%
We assume that $\varepsilon _{i}\perp \left( U_{i}^{\ast },V_{i}^{\ast
},W_{i},Z_{i}\right) $, i.e., the measurement error is classical.%

Note that the objects of researchers' interest do not change: the goal is to
estimate the partial effects~(\ref{eq:PE_Tobit})-(\ref{eq:PE_Probit}). The
structural endogeneity and measurement errors are difficulties that an
estimation procedure needs to overcome. In particular, note that we are
interested in estimation of the effect of $X_{i}^{\ast }$ on $Y_{i}$, and
not in the effect of the mismeasured $X_{i}$.\footnote{%
This is similar to the linear regression settings, where one would be
interested in the effect of $X_{i}^{\ast }$ on $Y_{i}$. The slope
coefficient in the OLS\ regression of $Y_{i}$ on $X_{i}$ is not the object
of interest because it is subject to the attenuation bias due to the EiV
(and also possibly due to the endogeneity of $X_{i}^{\ast }$).}

\bigskip

\section{Analysis of the Model\label{sec:Analysis}}

First, we use the model in equations (\ref{eq:dgp Y exog})-(\ref{eq:ME}) to
obtain the model in terms of the observable $X_{i}$. Since $X_{i}^{\ast
}=X_{i}-\varepsilon _{i}$, we can rewrite (\ref{eq:dgp Y exog}) as 
\begin{eqnarray*}
Y_{i} &=&m\left( \theta _{01}X_{i}^{\ast }+\theta _{02}^{\prime
}W_{i}+U_{i}^{\ast }\right) =m\left( \theta _{01}X_{i}+\theta _{02}^{\prime
}W_{i}-\theta _{01}\varepsilon _{i}+U_{i}^{\ast }\right) \\
&=&m\left( \theta _{01}X_{i}+\theta _{02}^{\prime }W_{i}+U_{i}\right) ,
\end{eqnarray*}%
where $U_{i}\equiv U_{i}^{\ast }-\theta _{01}\varepsilon _{i}$. Let $%
V_{i}\equiv V_{i}^{\ast }+\varepsilon _{i}$. The model in equations (\ref%
{eq:dgp Y exog})-(\ref{eq:ME}) can be written as%
\begin{eqnarray}
Y_{i} &=&m\left( \theta _{01}X_{i}+\theta _{02}^{\prime }W_{i}+U_{i}\right) ,
\label{eq:dgp Y wo stars gen} \\
X_{i} &=&\pi _{01}^{\prime }Z_{i}+\pi _{02}^{\prime }W_{i}+V_{i},
\label{eq:dgp X wo stars gen} \\
\left( 
\begin{array}{c}
U_{i} \\ 
V_{i}%
\end{array}%
\right) &\sim &N\left( \left( 
\begin{array}{c}
0 \\ 
0%
\end{array}%
\right) ,%
\begin{pmatrix}
\sigma _{U}^{2} & \sigma _{UV} \\ 
\sigma _{UV} & \sigma _{V}^{2}%
\end{pmatrix}%
\right) .  \label{eq:dgp UV wo stars gen}
\end{eqnarray}%
The definitions of $U_{i}$ and $V_{i}$ imply that%
\begin{equation}
\sigma _{U}^{2}=\sigma _{U^{\ast }}^{2}+\theta _{01}^{2}\sigma _{\varepsilon
}^{2},\quad \sigma _{V}^{2}=\sigma _{V^{\ast }}^{2}+\sigma _{\varepsilon
}^{2},\quad \sigma _{UV}=\sigma _{U^{\ast }V^{\ast }}-\theta _{01}\sigma
_{\varepsilon }^{2}.  \label{eq:sigma2 w and wo stars}
\end{equation}

Note that variables $X_{i},U_{i},V_{i}$ are the analogs of the true
variables $X_{i}^{\ast },U_{i}^{\ast },V_{i}^{\ast }$ that arise due to the
measurement errors $\varepsilon _{i}$. In the absence of measurement errors,
i.e., when $\varepsilon _{i}=0$, we have $X_{i}=X_{i}^{\ast }$, $%
U_{i}=U_{i}^{\ast }$, $V_{i}=V_{i}^{\ast }.$

The model in equations~(\ref{eq:dgp Y wo stars gen})-(\ref{eq:dgp UV wo
stars gen}) can be estimated by MLE or using the control variable two-step
approach described earlier. Specifically, both approaches will consistently
estimate parameters $\theta _{0}$ and the covariance matrix of the
unobservables in equation (\ref{eq:dgp UV wo stars gen}), i.e., $\sigma
_{U}^{2}$, $\sigma _{UV}$, and $\sigma _{V}^{2}$. Note that because the
model is nonlinear, the marginal effects~(\ref{eq:PE_Tobit})-(\ref%
{eq:PE_Probit}) depend not only on $\theta _{0}$ but also on $\sigma
_{U^{\ast }}^{2}$. Thus, even though the available data $\left(
Y_{i},X_{i},W_{i},Z_{i}\right) $ allows immediately estimating $\theta _{0}$%
, we cannot obtain the marginal effects because we do not know $\sigma
_{U^{\ast }}^{2}$. Naively using an estimate of $\sigma _{U}^{2}$ in place
of $\sigma _{U^{\ast }}^{2}$ would lead to a biased estimate of the partial
effects, since $\sigma _{U}^{2}\geq \sigma _{U^{\ast }}^{2}$, as implied by
equation~(\ref{eq:sigma2 w and wo stars}).

The problem with identifying $\sigma _{U^{\ast }}^{2}$ is that the data only
allows identification of the $3$ parameters $\sigma _{U}^{2}$, $\sigma _{UV}$%
, and $\sigma _{V}^{2}$. However, the distribution of the true $\left(
U_{i}^{\ast },V_{i}^{\ast },\varepsilon _{i}\right) $ is governed by $4$
parameters: $\sigma _{U^{\ast }}^{2}$, $\sigma _{U^{\ast }V^{\ast }}$, $%
\sigma _{V^{\ast }}^{2}$, and $\sigma _{\varepsilon }^{2}$. Thus, one cannot
uniquely determine these $4$ parameters from the $3$ equations~(\ref%
{eq:sigma2 w and wo stars}). In other words, models with different values of 
$\sigma _{\varepsilon }^{2}$ are observationally equivalent: they correspond
to identical distributions of the observables $\left(
Y_{i},X_{i},W_{i},Z_{i}\right) $ even though they imply different values of
true $\sigma _{U^{\ast }}^{2}$. Thus, one cannot uniquely determine (i.e.,
point-identify) $\sigma _{U^{\ast }}^{2}$ from the data $\left(
Y_{i},X_{i},W_{i},Z_{i}\right) $. Correspondingly, one cannot point-identify
the partial effects, which depend on $\sigma _{U^{\ast }}^{2}$.

Equations~(\ref{eq:sigma2 w and wo stars}) provide restrictions on $\sigma
_{U^{\ast }}^{2}$, which we will use to provide bounds on the possible
values of true $\sigma _{U^{\ast }}^{2}$, and hence on the values of the
partial effects.

\paragraph{Bounds on $\protect\sigma _{U^{\ast }}^{2}$ \label%
{sec:sig_Ustar_bounds}}

From equations~(\ref{eq:sigma2 w and wo stars}) the upper bound on $\sigma
_{U^{\ast }}^{2}$ is $\sigma _{U^{\ast }}^{2}\leq \sigma _{U}^{2}$. We now
obtain the lower bound on $\sigma _{U^{\ast }}^{2}$. In particular, we look
to find the smallest $\sigma _{U^{\ast }}^{2}$ that satisfies equations~(\ref%
{eq:sigma2 w and wo stars}), Cauchy-Schwarz inequality $\sigma _{U^{\ast
}V^{\ast }}^{2}\leq \sigma _{U^{\ast }}^{2}\sigma _{V^{\ast }}^{2}$, and the
non-negativity constraints $\sigma _{U^{\ast }}^{2}\geq 0$, $\sigma
_{V^{\ast }}^{2}\geq 0$, and $\sigma _{\varepsilon }^{2}\geq 0$. Let $\rho
_{UV}=\mathrm{corr}\left( U_{i},V_{i}\right) $.

\begin{proposition}
\label{prop:Ustar bounds}Suppose $\left\vert \rho _{UV}\right\vert <1$ in
model~(\ref{eq:dgp Y wo stars gen})-(\ref{eq:dgp UV wo stars gen}). Then%
\begin{equation*}
\sigma _{U^{\ast }}^{2}\in \left[ \underline{\sigma }_{U^{\ast }}^{2},\sigma
_{U}^{2}\right] ,
\end{equation*}%
where%
\begin{equation}
\underline{\sigma }_{U^{\ast }}^{2}\equiv \max \left\{ \frac{\left( \theta
_{01}\sigma _{UV}+\sigma _{U}^{2}\right) ^{2}}{\sigma _{V}^{2}\theta
_{01}^{2}+2\sigma _{UV}\theta _{01}+\sigma _{U}^{2}},\ \sigma
_{U}^{2}-\theta _{01}^{2}\sigma _{V}^{2}\right\} .  \label{eq:bound s2Us alt}
\end{equation}
\end{proposition}

Proposition \ref{prop:Ustar bounds} provides the bounds in terms of the
quantities that can be estimated using the data $\left(
Y_{i},X_{i},W_{i},Z_{i}\right) $. Condition $\left\vert \rho
_{UV}\right\vert <1$ guarantees that the denominator in the fraction above
is positive. The proof of Proposition \ref{prop:Ustar bounds} also provides
bounds on $\sigma _{U^{\ast }V^{\ast }}$ and $\sigma _{\varepsilon }^{2}$.

\paragraph{Correct Partial Effects\label{sec:correct PE}}

We now use the bounds on $\sigma _{U^{\ast }}^{2}$ from Proposition \textbf{%
\ref{prop:Ustar bounds}} to obtain the bounds on the partial effects, in
terms of the parameters that can be recovered from data.{}

{}For a given $\sigma _{U^{\ast }}^{2}$,
the partial effects for the $j^{th}$ covariate are, as in equations (\ref%
{eq:PE_Tobit})-(\ref{eq:PE_Probit}),%
\begin{equation}
PE_{j}^{\text{Tob}}\left( h,\sigma _{U^{\ast }}^{2}\right) =\Phi \left( 
\frac{\theta _{0}^{\prime }h}{\sigma _{U^{\ast }}}\right) \theta _{0j}\text{
\ \ and \ \ }PE_{j}^{\Pr }\left( h,\sigma _{U^{\ast }}^{2}\right) =\phi
\left( \frac{\theta _{0}^{\prime }h}{\sigma _{U^{\ast }}}\right) \frac{%
\theta _{0j}}{\sigma _{U^{\ast }}}.  \label{eq:PE_ProbitTobit_sigStar}
\end{equation}

The lower and upper bounds for partial effects for the $j^{th}$ covariate, $%
PE_{j}\left( h\right) $, are computed as 
\begin{equation}
\min_{\sigma _{U^{\ast }}^{2}\in \left[ \underline{\sigma }_{U^{\ast
}}^{2},\sigma _{U}^{2}\right] }PE_{j}\left( h,\underline{\sigma }_{U^{\ast
}}^{2}\right) \text{\quad and\quad }\max_{\sigma _{U^{\ast }}^{2}\in \left[ 
\underline{\sigma }_{U^{\ast }}^{2},\sigma _{U}^{2}\right] }PE_{j}\left(
h,\sigma _{U^{\ast }}^{2}\right) .  \label{eq:bounds_PE_generic}
\end{equation}

Function $PE_{j}^{\text{Tob}}\left( h,\sigma _{U^{\ast }}^{2}\right) $ in (%
\ref{eq:PE_ProbitTobit_sigStar}) is a monotone function of $\sigma _{U^{\ast
}}^{2}$, so the minimum and maximum in equation (\ref{eq:bounds_PE_generic})
are achieved on the boundaries of interval $\left[ \underline{\sigma }%
_{U^{\ast }}^{2},\sigma _{U}^{2}\right] $.

Function $PE_{j}^{\Pr }\left( h,\sigma _{U^{\ast }}^{2}\right) $ in
equation~(\ref{eq:PE_ProbitTobit_sigStar}) is not monotone in $\sigma
_{U^{\ast }}^{2}$, but the bounds in equation (\ref{eq:bounds_PE_generic})
for $PE_{j}^{\Pr }\left( h\right) $ can also be simplified. The minimum and
maximum over $\sigma _{U^{\ast }}^{2}\in \left[ \underline{\sigma }_{U^{\ast
}}^{2},\sigma _{U}^{2}\right] $ can be attained only at $\sigma _{U^{\ast
}}^{2}=\underline{\sigma }_{U^{\ast }}^{2}$, at $\sigma _{U^{\ast
}}^{2}=\sigma _{U}^{2}$, and, when $\left( \theta _{0}^{\prime }h\right)
^{2}\in \left[ \underline{\sigma }_{U^{\ast }}^{2},\sigma _{U}^{2}\right] $,
at $\sigma _{U^{\ast }}^{2}=\left( \theta _{0}^{\prime }h\right) ^{2}$.
Thus, one only needs to evaluate $PE_{j}^{\Pr }\left( h,\sigma _{U^{\ast
}}^{2}\right) $ at these $2$ or $3$ points to calculate the minimum and
maximum in equation~(\ref{eq:bounds_PE_generic}).

Since $\sigma _{U}\geq \sigma _{U^{\ast }}$, naively using $\sigma _{U}$
instead of $\sigma _{U^{\ast }}$ when calculating $PE_{j}^{\text{Tob}}\left(
h\right) $, would lead to attenuation bias when $\theta _{0}^{\prime }h>0$,
but would bias $PE_{j}^{\text{Tob}}\left( h\right) $ away from zero when $%
\theta _{0}^{\prime }h<0$, i.e., the EiV would make naive $PE_{j}^{\text{Tob}%
}\left( h,\sigma _{U}^{2}\right) $ over-estimate the partial effects $%
PE_{j}^{\text{Tob}}\left( h\right) $ in the latter case. Likewise, for the
probability, naively using $PE_{j}^{\Pr }\left( h,\sigma _{U}^{2}\right) $
can both under- and over-estimate the true partial effect $PE_{j}^{\Pr
}\left( h\right) $.

\paragraph{Estimation}

Using the standard two-step or MLE approaches described in Section \ref%
{sec:Model}, one obtains the estimates of $\theta _{0}$, $\sigma _{U}^{2}$, $%
\sigma _{V}^{2}$, and $\sigma _{UV}$ (and of their variance-covariance
matrix for inference). Then, from equation (\ref{eq:bound s2Us alt}) one
obtains the estimate of $\underline{\sigma }_{U^{\ast }}^{2}$.

For a given value of $\sigma _{U^{\ast }}^{2}$, the estimated partial
effects would be%
\begin{equation}
\widehat{PE}_{j}^{\text{Tob}}\left( h,\sigma _{U^{\ast }}^{2}\right) =\Phi
\left( \frac{\widehat{\theta }^{\prime }h}{\sigma _{U^{\ast }}}\right) 
\widehat{\theta }_{j}\text{ \ \ and \ \ }\widehat{PE}_{j}^{\Pr }\left(
h,\sigma _{U^{\ast }}^{2}\right) =\phi \left( \frac{\widehat{\theta }%
^{\prime }h}{\sigma _{U^{\ast }}}\right) \frac{\widehat{\theta }_{j}}{\sigma
_{U^{\ast }}}.  \label{eq:hat PE-ProbitTobit(sigUs2)}
\end{equation}%
 Then, the estimated bounds on $%
PE_{j}\left( h\right) $ are %
\begin{equation}
\min_{v\in \left[ \widehat{\underline{\sigma }}_{U^{\ast }}^{2},\widehat{%
\sigma }_{U}^{2}\right] }\widehat{PE}_{j}\left( h,v\right) \text{\quad
and\quad }\max_{v\in \left[ \widehat{\underline{\sigma }}_{U^{\ast }}^{2},%
\widehat{\sigma }_{U}^{2}\right] }\widehat{PE}_{j}\left( h,v\right) ,
\label{eq:bounds_hat_PE_generic}
\end{equation}%
where the minimum and maximum are easily computed using univariate numerical
optimization. For the partial effects in equation (\ref{eq:hat
PE-ProbitTobit(sigUs2)}), these extrema can also be computed as described
under equation (\ref{eq:bounds_PE_generic}).

For example, one often considers the partial effects at the mean values of
covariates taking $h=(\overline{X},\overline{W}^{\prime })^{\prime }$, where 
$\overline{X}$ and $\overline{W}$ are the sample averages.  Note
that $\mathrm{E}\left[ X_{i}\right] =\mathrm{E}\left[ X_{i}^{\ast }\right] $.

\paragraph{Inference}

To provide a simple method for inference about the partial effects, we adopt
a Bonferroni approach (e.g., \citealp{McCloskey2017JoE}). This approach
allows us to avoid computational challenges that often arise in the context
of subvector inference in partially identified models. The construction of a 
$1-\alpha $ confidence interval for a partial effect $PE_{j}(h)$ proceeds in
two steps:

\begin{enumerate}
\item Pick $\alpha _{1}\in (0,\alpha )$ and construct $CI_{1-\alpha
_{1}}^{\sigma _{U^{\ast }}^{2}}$, a $1-\alpha _{1}$ confidence interval for $%
\sigma _{U^{\ast }}^{2}$, based on the bounds provided in Proposition~%
\textbf{\ref{prop:Ustar bounds}}.

\item Construct a $1-\alpha $ confidence interval for $PE_{j}(h)$ as the
union $CI_{1-\alpha }^{PE_{j}(h)}=\bigcup_{\sigma _{U^{\ast }}^{2}\in
CI_{1-\alpha _{1}}^{\sigma _{U^{\ast }}^{2}}}CI_{1-\left( \alpha -\alpha
_{1}\right) }^{PE_{j}(h)}\left( \sigma _{U^{\ast }}^{2}\right) $, where $%
CI_{1-\left( \alpha -\alpha _{1}\right) }^{PE_{j}(h)}\left( \sigma _{U^{\ast
}}^{2}\right) $ is a standard $1-(\alpha -\alpha _{1})$ confidence interval
for $PE_{j}(h)$ based on $\widehat{PE}_{{j}}\left( h,\sigma _{U^{\ast
}}^{2}\right) $ in equation~\eqref{eq:hat PE-ProbitTobit(sigUs2)} for a
given $\sigma _{U^{\ast }}^{2}$.
\end{enumerate}

We now provide the implementation details for each step.

\smallskip

\noindent \textit{Step 1}. The confidence interval for $\sigma _{U^{\ast
}}^{2}$ is constructed based on the bounds given in Proposition~\textbf{\ref%
{prop:Ustar bounds}}. As the upper bound, we take $\hat{\sigma}%
_{U}^{2}+z_{1-\alpha _{1}/2}\times s_{\hat{\sigma}_{U}^{2}}$, where $s_{\hat{%
\sigma}_{U}^{2}}$ is the standard error of $\hat{\sigma}_{U}^{2}$, and $%
z_{1-\alpha _{1}/2}$ is the $1-\alpha _{1}/2$ quantile of the standard
normal distribution. The lower bound is based on $\underline{\hat{\sigma}}%
_{U^{\ast }}^{2}=\max \{\hat{\xi}_{1},\hat{\xi}_{2}\}$, where $\hat{\xi}_{1}$
and $\hat{\xi}_{2}$ are the plug-in estimators of the two terms on the right
hand side of equation~\eqref{eq:bound s2Us alt}. Note that $\hat{\xi}_{1}$
and $\hat{\xi}_{2}$ are (generally) jointly asymptotically normal and their
asymptotic variance-covariance matrix can be computed using the delta
method. Then, as the lower bound of $CI_{1-\alpha _{1}}^{\sigma _{U^{\ast
}}^{2}}$, we take $\max \{\hat{\xi}_{1}-c_{1-\alpha _{1}/2}\times s_{\hat{\xi%
}_{1}},\hat{\xi}_{2}-c_{1-\alpha _{1}/2}\times s_{\hat{\xi}_{2}}\}$. Here $%
s_{\hat{\xi}_{1}}$ and $s_{\hat{\xi}_{2}}$ are the standard errors of $\hat{%
\xi}_{1}$ and $\hat{\xi}_{2}$, and $c_{1-\alpha _{1}/2}$ is the $1-\alpha
_{1}/2$ quantile of $\max \{\eta _{1},\eta _{2}\}$, where $\left( \eta
_{1},\eta _{2}\right) $ are jointly normal with unit variances and
correlation $\hat{\rho}_{\hat{\xi}_{1},\hat{\xi}_{2}}$, and $\hat{\rho}_{%
\hat{\xi}_{1},\hat{\xi}_{2}}$ is an estimator of the correlation between $%
\hat{\xi}_{1}$ and $\hat{\xi}_{2}$ (e.g., see \citealp{RomanoWolf2005Ecta}).
By a standard argument, the confidence interval for $\sigma _{U^{\ast }}^{2}$
given by 
\begin{equation*}
CI_{1-\alpha _{1}}^{\sigma _{U^{\ast }}^{2}}=\left[ \max \left\{ \hat{\xi}%
_{1}-c_{1-\alpha _{1}/2}\times s_{\hat{\xi}_{1}},\hat{\xi}_{2}-c_{1-\alpha
_{1}/2}\times s_{\hat{\xi}_{2}}\right\} ,\hat{\sigma}_{U}^{2}+z_{1-\alpha
_{1}/2}\times s_{\hat{\sigma}_{U}^{2}}\right]
\end{equation*}%
has asymptotic coverage at least $1-\alpha _{1}$ for the true $\sigma
_{U^{\ast }}^{2}$. In the numerical illustrations we take $\alpha
_{1}=\alpha /10$.

\smallskip

\noindent \textit{Step 2}. First, the standard $CI_{1-\left( \alpha -\alpha
_{1}\right) }^{PE_{j}(h)}\left( \sigma _{U^{\ast }}^{2}\right) $ is $\left[
l_{1-\left( \alpha -\alpha _{1}\right) }^{PE_{j}(h)}\left( \sigma _{U^{\ast
}}^{2}\right) ,u_{1-\left( \alpha -\alpha _{1}\right) }^{PE_{j}(h)}\left(
\sigma _{U^{\ast }}^{2}\right) \right]$ constructed by adding and
subtracting $z_{1-(\alpha -\alpha _{1})/2}\times s_{\widehat{PE}%
_{j}(h,\sigma _{U^{\ast }}^{2})}$ from $\widehat{PE}_{j}(h,\sigma _{U^{\ast
}}^{2})$. The standard error $s_{\widehat{PE}_{j}(h,\sigma _{U^{\ast
}}^{2})} $ of $\widehat{PE}_{j}(h,\sigma _{U^{\ast }}^{2})$ can be computed
using the delta method. Then we can construct $CI_{1-\alpha }^{PE_{j}(h)}$
as 
\begin{equation*}
CI_{1-\alpha }^{PE_{j}(h)} = \left[\min_{\sigma _{U^{\ast }}^{2}\in
CI_{1-\alpha _{1}}^{\sigma _{U^{\ast }}^{2}}}l_{1-\left( \alpha -\alpha
_{1}\right) }^{PE_{j}(h)}\left( \sigma _{U^{\ast }}^{2}\right) ,
\max_{\sigma _{U^{\ast }}^{2}\in CI_{1-\alpha _{1}}^{\sigma _{U^{\ast
}}^{2}}}u_{1-\left( \alpha -\alpha _{1}\right) }^{PE_{j}(h)}\left( \sigma
_{U^{\ast }}^{2}\right)\right] ,
\end{equation*}%
where the minimum and maximum are easily calculated using univariate
numerical optimization over $\sigma _{U^{\ast }}^{2}$. By the standard
Bonferroni argument, the confidence interval $CI_{1-\alpha }^{PE_{j}(h)}$
has asymptotic coverage of at least $1-\alpha $ for the true partial effect $%
PE_{j}(h)$.

\smallskip

The constructed confidence interval is asymptotically valid as long as (i)
the first step confidence interval $CI_{1-\alpha _{1}}^{\sigma _{U^{\ast
}}^{2}}$ covers the true $\sigma _{U^{\ast }}^{2}$ with probability at least 
$1-\alpha _{1}$ asymptotically, and (ii) the delta method applies to $%
\widehat{PE}_{j}(h,\sigma _{U^{\ast }}^{2})$ for the true $\sigma _{U^{\ast
}}^{2}$. Both conditions are satisfied provided that the true $\sigma
_{U^{\ast }}^{2}$ is bounded away from zero. Note that in this case $%
CI_{1-\alpha _{1}}^{\sigma _{U^{\ast }}^{2}}$ is valid even if $\theta
_{01}\sigma _{UV}+\sigma _{U}^{2}$ is equal to (or local to) zero, which
implies that $CI_{1-\alpha }^{PE_{j}(h)}$ is also valid.

\section{Extensions\label{sec:Extensions}}

\subsection{Probit\label{sec:Probit}}

IV-Probit is the same as IV-Tobit except $m\left( s\right) =1\left\{
s>0\right\} $ in equation (\ref{eq:dgp Y exog}). Since Probit is a binary
outcome model, in equation~(\ref{eq:dgp UV wo stars gen}) one imposes the
standard normalization $\sigma _{U}=1$. For Probit, we are interested in the
partial effects of covariates on the probability of $Y_{i}=1$, which are
given by $PE_{j}^{\Pr }\left( h\right) $ in equation (\ref{eq:PE_Probit}).\
Similarly to the IV-Tobit model, the IV-Probit model can be estimated by MLE
or by the two-step approach identical to the one described in Section~\ref%
{sec:Model}, except the second step uses the standard Probit estimator in
place of the Tobit estimator. Then the bounds on $PE_{j}^{\Pr }\left(
h\right) $ are estimated as in equation (\ref{eq:bounds_hat_PE_generic}).%
{} Confidence intervals for $PE_{j}^{\Pr }(h)$ can be computed exactly
as described above.

\subsection{Average Partial Effects and other Other Counterfactuals \label%
{sec:APEs}}

In addition to the partial effects at a given $h$, researchers are often
interested in the Average Partial Effects%
\begin{equation}
APE_{j}^{\text{Tob}}\equiv \mathrm{E}\left[ PE_{j}^{\text{Tob}}\left(
H_{i}^{\ast }\right) \right] \text{ and }APE_{j}^{\Pr }\equiv \mathrm{E}%
\left[ PE_{j}^{\Pr }\left( H_{i}^{\ast }\right) \right] ,
\label{eq:APEs def}
\end{equation}%
which are the partial effects $PE_{j}\left( h\right) $ averaged with respect
to the distribution of {}$H_{i}^{\ast }=\left(
X_{i}^{\ast },W_{i}^{\prime }\right) ^{\prime }$. Define%
\begin{equation*}
APE_{j}^{\text{Tob}}\left( \sigma _{U^{\ast }}^{2}\right) \equiv \mathrm{E}%
\left[ PE_{j}^{\text{Tob}}\left( H_{i}^{\ast },\sigma _{U^{\ast
}}^{2}\right) \right] \text{ and }APE_{j}^{\Pr }\left( \sigma _{U^{\ast
}}^{2}\right) \equiv \mathrm{E}\left[ PE_{j}^{\Pr }\left( H_{i}^{\ast
},\sigma _{U^{\ast }}^{2}\right) \right] .
\end{equation*}

Note that the distribution of $X_{i}^{\ast }$ is not directly observable due
to the EiV. Averaging $PE_{j}\left( h\right) $ with respect to the
distribution of the observed $H_{i}=\left( X_{i},W_{i}^{\prime }\right)
^{\prime }$ would result in biased estimators of the APEs. To account for
this, in the Appendix we show that these APEs can be calculated as%
\begin{eqnarray}
APE_{j}^{\text{Tob}}\left( \sigma _{U^{\ast }}^{2}\right) &=&\mathrm{E}\left[
\Phi \left( \frac{\theta _{01}\pi _{01}^{\prime }Z_{i}+\left( \theta
_{01}\pi _{02}+\theta _{02}\right) ^{\prime }W_{i}}{\sqrt{2\sigma _{U^{\ast
}}^{2}-\sigma _{U}^{2}+\theta _{01}^{2}\sigma _{V}^{2}}}\right) \right]
\theta _{0j},  \label{eq:true APE Tobit} \\
APE_{j}^{\Pr }\left( \sigma _{U^{\ast }}^{2}\right) &=&\mathrm{E}\left[ \phi
\left( \frac{\theta _{01}\pi _{01}^{\prime }Z_{i}+\left( \theta _{01}\pi
_{02}+\theta _{02}\right) ^{\prime }W_{i}}{\sqrt{2\sigma _{U^{\ast
}}^{2}-\sigma _{U}^{2}+\theta _{01}^{2}\sigma _{V}^{2}}}\right) \right] 
\frac{\theta _{0j}}{\sqrt{2\sigma _{U^{\ast }}^{2}-\sigma _{U}^{2}+\theta
_{01}^{2}\sigma _{V}^{2}}}.  \label{eq:true APE Probit}
\end{eqnarray}%
Hence, for any given value of $\sigma _{U^{\ast }}^{2}$, these APEs can be
estimated by%
\begin{eqnarray*}
\widehat{APE}_{j}^{\text{Tob}}\left( \sigma _{U^{\ast }}^{2}\right) &=&\frac{%
1}{n}\sum_{i=1}^{n}\Phi \left( \frac{\hat{\theta}_{1}\hat{\pi}_{1}^{\prime
}Z_{i}+\left( \hat{\theta}_{1}\hat{\pi}_{2}+\hat{\theta}_{2}\right) ^{\prime
}W_{i}}{\sqrt{2\sigma _{U^{\ast }}^{2}-\hat{\sigma}_{U}^{2}+\hat{\theta}%
_{1}^{2}\hat{\sigma}_{V}^{2}}}\right) \hat{\theta}_{j}, \\
\widehat{APE}_{j}^{\Pr }\left( \sigma _{U^{\ast }}^{2}\right) &=&\frac{1}{n}%
\sum_{i=1}^{n}\phi \left( \frac{\hat{\theta}_{1}\hat{\pi}_{1}^{\prime
}Z_{i}+\left( \hat{\theta}_{1}\hat{\pi}_{2}+\hat{\theta}_{2}\right) ^{\prime
}W_{i}}{\sqrt{2\sigma _{U^{\ast }}^{2}-\hat{\sigma}_{U}^{2}+\hat{\theta}%
_{1}^{2}\hat{\sigma}_{V}^{2}}}\right) \frac{\hat{\theta}_{j}}{\sqrt{2\sigma
_{U^{\ast }}^{2}-\hat{\sigma}_{U}^{2}+\hat{\theta}_{1}^{2}\hat{\sigma}%
_{V}^{2}}}.
\end{eqnarray*}%
Finally, the estimated bounds on the APEs are obtained by finding the
minimum and maximum over $\sigma _{U^{\ast }}^{2}\in \left[ \hat{\underline{%
\sigma }}_{U^{\ast }}^{2},\hat{\sigma}_{U}^{2}\right] $. These can be easily
computed numerically, since $\widehat{APE}_{j}\left( v\right) $ are smooth
functions of a scalar argument $v$. Our two-step approach to inference also
applies to the APEs with a minimal modification. The only difference is that
in Step 2 the construction of the standard error $s_{\widehat{APE}%
_{j}}\left( \sigma _{U^{\ast }}^{2}\right) $ as usual needs to account for
the sampling variability in both the parameter estimators and the data
entering the expressions for the APEs directly.

It is also straightforward to apply our analysis to other counterfactuals,
including partial effects and APEs of discrete covariates, as well as to the
ordered Probit and two-sided Tobit models. Proposition~\ref{prop:Ustar
bounds} and the bounds on $\sigma _{U^{\ast }}^{2}$ remain the same, and
hence the estimation and inference procedures remain unchanged, except for
different formulas in equations~\eqref{eq:PE_ProbitTobit_sigStar}-%
\eqref{eq:bounds_hat_PE_generic} corresponding to the counterfactuals of
interest.

\section{Numerical Illustration\label{sec:MC}}

We simulate a Tobit model with endogenous and mismeasured $X_{i}^{\ast }$,
as in equations (\ref{eq:dgp Y exog})-(\ref{eq:ME}), with $W_{i}=1$, $%
Z_{i}\sim N\left( 0,1\right) $, $\left( \theta _{01},\theta _{02},\sigma
_{V^{\ast }},\sigma _{U^{\ast }},\sigma _{\varepsilon },\pi _{01},\pi
_{02}\right) =\left( 2,1,1,1,1,1,0\right) $, and $n=200$. {}Figure~\ref%
{fig:MC} plots the results for the Partial Effects (PEs) of $X_{i}^{\ast }$
at the population mean values of the covariates.

We consider a range of designs corresponding to the true values of $\rho
_{U^{\ast }V^{\ast }}\in \left[ -0.95,0.95\right] $ on the horizontal axis.
For each $\rho _{U^{\ast }V^{\ast }}$, the figure shows the true PE
(\textquotedblleft true\textquotedblright ), the true (population) bounds
for the PE obtained using Proposition~\ref{prop:Ustar bounds}
(\textquotedblleft true bounds\textquotedblright ), as well as the medians
over the Monte Carlo replications of the estimated lower and upper bounds on
the PE (\textquotedblleft LB\textquotedblright\ and \textquotedblleft
UB\textquotedblright ) and the corresponding $95\%$ confidence intervals
(\textquotedblleft CI\textquotedblright ) based on the two-step IV-Tobit
estimator. The true bounds for the PE are calculated using the point
identified parameters $\theta _{0}$, $\sigma _{U}^{2}$, $\sigma _{UV}$, and $%
\sigma _{V}^{2}$, see equations (\ref{eq:dgp Y wo stars gen})-(\ref{eq:dgp
UV wo stars gen}). For comparison, we also include the results for the PE
calculated using the standard naive IV-Tobit estimator (\textquotedblleft
naive\textquotedblright ) and the corresponding confidence intervals
(\textquotedblleft CI naive\textquotedblright ). The \textquotedblleft
naive\textquotedblright\ estimators of the partial effects are $\widehat{PE}%
_{1}^{\text{Tob}}\left( h,\widehat{\sigma }_{U}^{2}\right) $ and $\widehat{PE%
}_{1}^{\Pr }\left( h,\widehat{\sigma }_{U}^{2}\right) $, i.e., they replace $%
\sigma _{U^{\ast }}^{2}$ with $\widehat{\sigma }_{U}^{2}$ in equation (\ref%
{eq:hat PE-ProbitTobit(sigUs2)}).

\begin{figure}[!h]
\centering%
\includegraphics[trim = {0.5cm 0.7cm 0.5cm
0.2cm},clip,width=\textwidth,keepaspectratio]{figure1.pdf}
\caption{Simulation results for partial effects on: (a) expectations, $%
PE_{h_{1}}^{\text{Tob}}\left( h\right)$, and (b) probability, $PE_{1}^{\text{%
Pr}}\left( h\right) $. Values of $\protect\rho_{U^{\ast }V^{\ast }}$ are on
the horizontal axis.}
\label{fig:MC}
\end{figure}

Figure~\ref{fig:MC} (a) shows the bounds on $PE_{h_{1}}^{\text{Tob}}\left(
h\right) $, while Figure~\ref{fig:MC} (b) considers $PE_{1}^{\text{Pr}%
}\left( h\right) $. As expected, the true PE is between the lower and upper
bounds for all values of $\rho _{U^{\ast }V^{\ast }}$. By construction, the
\textquotedblleft naive\textquotedblright\ \ IV-Tobit estimator of $PE_{1}^{%
\text{Tob}}\left( h\right) $ coincides with one of the bounds. In both
panels, the \textquotedblleft naive\textquotedblright\ estimates are below
the true values for every $\rho _{U^{\ast }V^{\ast }}$, and the
\textquotedblleft naive\textquotedblright\ IV-Tobit confidence intervals do
not include the true PE. In this design, the identified set and the
confidence intervals for the true partial effects are much wider than those
of the \textquotedblleft naive\textquotedblright\ estimator. The relative
width of the identified set and the confidence intervals depends on the
specific parameter values. In contrast to these simulations, in the example
of the next section, the identified set is very narrow and the confidence
intervals for the partial effects have width similar to those of the naive
estimator.

To gain some intuition about the shape of the bounds in Figure~\ref{fig:MC},
notice that when $\theta _{01}>0$, measurement error in $X_{i}^{\ast }$
introduces a negative correlation between $X_{i}$ and $U_{i}$. Thus,
observing $\rho _{UV}>0$ is only consistent with $\rho _{U^{\ast }V^{\ast
}}>0$, but not with $\rho _{U^{\ast }V^{\ast }}\leq 0$, i.e., implies
positive correlation due to the structural endogeneity. On the other hand,
observing $\rho _{UV}<0$ can be explained both by the effect of EiV combined
with $\rho _{U^{\ast }V^{\ast }}\gtrless 0$, and by $\rho _{U^{\ast }V^{\ast
}}<0$ without any EiV. Thus, when $\rho _{UV}<0$ it is harder to disentangle
structural endogeneity and EiV. As a consequence, the true and the estimated
correct bounds on the PEs in Figure~\ref{fig:MC} are wider for the negative
values of $\rho _{U^{\ast }V^{\ast }}$.

\section{Empirical Illustration\label{sec:Empirical}}

We illustrate the proposed methods in the estimation of the Tobit and Probit
models for women's labor force participation. We use the classic data set
from \cite{Mroz1987Ecta} who estimates Tobit and related models to explain
married women's hours of work. The data contains 753 women, 428 of which
report working non-zero hours. For Tobit, the dependent variable is the
number of hours worked (\textit{hours}), and for Probit, the dependent
variable is working at some point during the year ($1\{hours>0\}$). In both
models, the covariates are age, education, experience, experience squared,
nonwife income in thousands (\textit{nwifeinc}), number of children less
than six years of age, number of children between 6 and 18 inclusive, and an
intercept. Husband's years of schooling, \textit{huseduc}, are used as an
instrument for \textit{nwifeinc}. This specification is used in, e.g., \cite%
{Wooldridge2010Book}. Nonwife income could be correlated with the unobserved
characteristics (structural endogeneity), and income variables are also
known to be frequently mismeasured.

\begin{table}[!h]
\rowcolors{1}{white}{myEvenRowsColor}
  \begin{center}
  \begin{tabular}{|c|c|cc|cc|}
    \hline\hline
   &  Tobit & IV-Tobit & CI for IV-Tobit & [LB, UB] & CI \\ \hline
    nwifeinc &  -5.33 & -19.0 & [\,-39.6,\,1.68\,] & [\,-19.1,\,-19.0\,] & [\,-41.6,\,2.44\,]\\ 
    educ &  48.7 & 70.3 & [\,29.0,\,112\,] & [\,70.3,\,70.8\,] & [\,26.9,\,117\,]\\ 
    exper &  79.5 & 74.9 & [\,51.6,\,98.2\,] & [\,74.9,\,75.4\,] & [\,50.3,\,102\,]\\ 
    exper2 &  -1.13 & -1.14 & [\,-1.82,\,-0.468\,] & [\,-1.15,\,-1.14\,] & [\,-1.89,\,-0.444\,]\\ 
    age &  -32.9 & -28.2 & [\,-39.3,\,-17.2\,] & [\,-28.4,\,-28.2\,] & [\,-40.6,\,-16.8\,]\\ 
  \hline\hline
  \end{tabular}
  \caption{Tobit. Partial Effects on Expectation.}
  \label{table:Emp-Tobit-PEM-E-w-CI}
  \end{center}
\end{table}

\begin{table}[!h]
\rowcolors{1}{white}{myEvenRowsColor}
  \begin{center}
  \begin{tabular}{|c|c|cc|cc|}
    \hline\hline
   &  Tobit & IV-Tobit & CI for IV-Tobit & [LB, UB] & CI \\ \hline
    nwifeinc &  -0.303 & -1.06 & [\,-2.16,\,0.043\,] & [\,-1.10,\,-1.06\,] & [\,-2.65,\,0.157\,]\\ 
    educ &  2.77 & 3.92 & [\,1.75,\,6.10\,] & [\,3.92,\,4.08\,] & [\,1.33,\,7.48\,]\\ 
    exper &  4.52 & 4.18 & [\,2.77,\,5.59\,] & [\,4.18,\,4.34\,] & [\,2.51,\,6.51\,]\\ 
    exper2 &  -0.064 & -0.064 & [\,-0.102,\,-0.026\,] & [\,-0.066,\,-0.064\,] & [\,-0.121,\,-0.022\,]\\ 
    age &  -1.87 & -1.58 & [\,-2.26,\,-0.890\,] & [\,-1.64,\,-1.58\,] & [\,-2.60,\,-0.834\,]\\ 
  \hline\hline
  \end{tabular}
  \caption{Tobit. Partial Effects on Probability. All numbers are multiplied by 100.}
  \label{table:Emp-Tobit-PEM-P-w-CI}
  \end{center}
\end{table}


\begin{table}[!h]
\rowcolors{1}{white}{myEvenRowsColor}
  \begin{center}
  \begin{tabular}{|c|c|cc|cc|}
    \hline\hline
   &  Probit & IV-Probit & CI for IV-Probit & [LB, UB] & CI \\ \hline
    nwifeinc &  -0.470 & -1.39 & [\,-2.67,\,-0.104\,] & [\,-1.49,\,-1.39\,] & [\,-3.29,\,0.079\,]\\ 
    educ &  5.11 & 6.41 & [\,3.96,\,8.86\,] & [\,6.41,\,6.87\,] & [\,2.98,\,10.8\,]\\ 
    exper &  4.82 & 4.38 & [\,2.68,\,6.08\,] & [\,4.38,\,4.70\,] & [\,2.49,\,6.82\,]\\ 
    exper2 &  -0.074 & -0.073 & [\,-0.118,\,-0.028\,] & [\,-0.079,\,-0.073\,] & [\,-0.137,\,-0.024\,]\\ 
    age &  -2.06 & -1.69 & [\,-2.58,\,-0.804\,] & [\,-1.81,\,-1.69\,] & [\,-2.87,\,-0.784\,]\\ 
  \hline\hline
  \end{tabular}
  \caption{Probit. Partial Effects on Probability. All numbers are multiplied by 100.}
  \label{table:Emp-Probit-PEM-P-w-CI}
  \end{center}
\end{table}


Tables \ref{table:Emp-Tobit-PEM-E-w-CI}-\ref{table:Emp-Tobit-PEM-P-w-CI}
contain the results on partial effects for Tobit. Table \ref%
{table:Emp-Tobit-PEM-E-w-CI} contains the results on partial effects on
expectation, $PE^{\text{Tob}}$, while Table \ref{table:Emp-Tobit-PEM-P-w-CI}
contains partial effects on probability, $PE^{\text{Pr}}$. All partial
effects are evaluated at the mean values of covariates. In both tables, the
first column (\textquotedblleft Tobit") provides the partial effects for
different covariates in the standard Tobit MLE where all covariates are
assumed to be exogenous. 
The remaining columns are based on the two-step IV-Tobit estimator, where 
\textit{huseduc} is used to instrument for the endogenous \textit{nwifeinc.}
The second column (\textquotedblleft IV-Tobit") contains the naive
estimators of the partial effects, followed by the $95\%$ confidence
intervals (column \textquotedblleft CI for IV-Tobit"). Column
\textquotedblleft \lbrack LB, UB]" provides the proposed estimated bounds
for the partial effects that account for both types of endogeneity. The last
column contains the corresponding confidence intervals for the partial
effects.

In both Tables \ref{table:Emp-Tobit-PEM-E-w-CI} and \ref%
{table:Emp-Tobit-PEM-P-w-CI}, we observe that the confidence intervals for
the correct partial effects at the mean are only slightly wider than the
naive ones of IV-Tobit. In particular, using the correct inference approach
does not change any of the conclusions about the effects of the variables
being statistically significant.

Table \ref{table:Emp-Probit-PEM-P-w-CI} contains the corresponding results
for Probit. Again, the confidence intervals for the correct partial effects
are not much wider than the naive ones. Unlike the Tobit case, not all
conclusions about the statistical significance of the partial effects are
preserved, as the partial effect of nonwife income %
becomes statistically insignificant. Hence, if we base our analysis only on
the binary outcome model, properly accounting for the roles of measurement
error and structural endogeneity can reverse the conclusion that the partial
effect of nonwife income %
is statistically significantly different from zero.

\bibliographystyle{ecta}
\bibliography{ME_Probit}

\section*{Appendix}

\singlespacing\CMDsmall

\noindent \textbf{Proof of Proposition~\ref{prop:Ustar bounds}}$.$ Note that 
$\sigma _{U^{\ast }V^{\ast }}^{2}\leq \sigma _{U^{\ast }}^{2}\sigma
_{V^{\ast }}^{2}$ combined with equation~(\ref{eq:sigma2 w and wo stars})
implies%
\begin{eqnarray*}
0 &\leq &\sigma _{U^{\ast }}^{2}\sigma _{V^{\ast }}^{2}-\sigma _{U^{\ast
}V^{\ast }}^{2}=\left( \sigma _{U}^{2}-\theta _{01}^{2}\sigma _{\varepsilon
}^{2}\right) \left( \sigma _{V}^{2}-\sigma _{\varepsilon }^{2}\right)
-\left( \sigma _{UV}+\theta _{01}\sigma _{\varepsilon }^{2}\right) ^{2}%
{}
\\
&=&\sigma _{U}^{2}\sigma _{V}^{2}-\sigma _{UV}^{2}-\sigma _{\varepsilon
}^{2}\left( \sigma _{V}^{2}\theta _{01}^{2}+2\sigma _{UV}\theta _{01}+\sigma
_{U}^{2}\right) ,\text{ and hence} \\
0 &\leq &\sigma _{\varepsilon }^{2}\leq \frac{\sigma _{U}^{2}\sigma
_{V}^{2}-\sigma _{UV}^{2}}{\sigma _{V}^{2}\theta _{01}^{2}+2\sigma
_{UV}\theta _{01}+\sigma _{U}^{2}}.
\end{eqnarray*}%
Since $\left\vert \rho _{UV}\right\vert <1$, the denominator in this
fraction is positive. Since $\sigma _{\varepsilon }^{2}\leq \sigma _{V}^{2}$%
, let%
\begin{equation*}
\overline{\sigma }_{\varepsilon }^{2}\equiv \min \left\{ \frac{\sigma
_{U}^{2}\sigma _{V}^{2}-\sigma _{UV}^{2}}{\sigma _{V}^{2}\theta
_{01}^{2}+2\sigma _{UV}\theta _{01}+\sigma _{U}^{2}},\sigma _{V}^{2}\right\}
.
\end{equation*}%
Then%
\begin{equation}
\underline{\sigma }_{U^{\ast }}^{2}\leq \sigma _{U^{\ast }}^{2}\leq \sigma
_{U}^{2},\qquad \underline{\sigma }_{U^{\ast }}^{2}\equiv \max \left\{
\sigma _{U}^{2}-\theta _{01}^{2}\overline{\sigma }_{\varepsilon
}^{2},0\right\} .  \label{eq:bound s2Us}
\end{equation}%
Note that%
\begin{equation*}
\sigma _{U}^{2}-\theta _{01}^{2}\frac{\sigma _{U}^{2}\sigma _{V}^{2}-\sigma
_{UV}^{2}}{\sigma _{V}^{2}\theta _{01}^{2}+2\sigma _{UV}\theta _{01}+\sigma
_{U}^{2}}=\frac{\sigma _{U}^{4}+2\sigma _{U}^{2}\sigma _{UV}\theta
_{01}+\theta _{01}^{2}\sigma _{UV}^{2}}{\sigma _{V}^{2}\theta
_{01}^{2}+2\sigma _{UV}\theta _{01}+\sigma _{U}^{2}}=\frac{\left( \theta
_{01}\sigma _{UV}+\sigma _{U}^{2}\right) ^{2}}{\sigma _{V}^{2}\theta
_{01}^{2}+2\sigma _{UV}\theta _{01}+\sigma _{U}^{2}}.
\end{equation*}%
Thus, $\underline{\sigma }_{U^{\ast }}^{2}$ in equation~(\ref{eq:bound s2Us}%
) can be equivalently written as%
\begin{equation*}
\underline{\sigma }_{U^{\ast }}^{2}\equiv \max \left\{ \frac{\left( \theta
_{01}\sigma _{UV}+\sigma _{U}^{2}\right) ^{2}}{\sigma _{V}^{2}\theta
_{01}^{2}+2\sigma _{UV}\theta _{01}+\sigma _{U}^{2}},\sigma _{U}^{2}-\theta
_{01}^{2}\sigma _{V}^{2}\right\} ,
\end{equation*}%
where the fraction is always non-negative, ensuring that $\underline{\sigma }%
_{U^{\ast }}^{2}\geq 0$. Finally, equation~(\ref{eq:sigma2 w and wo stars})
implies that $\sigma _{U^{\ast }V^{\ast }}$ is bounded between $\sigma _{UV}$
and $\sigma _{UV}+\theta _{01}\overline{\sigma }_{\varepsilon }^{2}$.\hfill $%
\blacksquare $

\bigskip

\noindent \textbf{Derivation of equation~(\ref{eq:true APE Tobit})}: 
\begin{align*}
& APE_{j}^{\text{Tob}}\left( \sigma _{U^{\ast }}^{2}\right) =\mathrm{E}\left[
\Phi \left( \frac{\theta _{0}^{\prime }H_{i}^{\ast }}{\sigma _{U^{\ast }}}%
\right) \theta _{0j}\right] =\mathrm{E}\left[ \Phi \left( \frac{\theta
_{01}X_{i}^{\ast }+\theta _{02}^{\prime }W_{i}}{\sigma _{U^{\ast }}}\right) %
\right] \theta _{0j}{} \\
& =\mathrm{E}\left[ \Phi \left( \frac{\theta _{01}\pi _{01}^{\prime
}Z_{i}+\left( \theta _{01}\pi _{02}+\theta _{02}\right) ^{\prime }W_{i}}{%
\sigma _{U^{\ast }}}+\frac{\theta _{01}V_{i}^{\ast }}{\sigma _{U^{\ast }}}%
\right) \right] \theta _{0j}=\mathrm{E}\left[ \Phi \left( \frac{\theta
_{01}\pi _{01}^{\prime }Z_{i}+\left( \theta _{01}\pi _{02}+\theta
_{02}\right) ^{\prime }W_{i}}{\sigma _{U^{\ast }}\sqrt{1+\theta
_{01}^{2}\sigma _{V^{\ast }}^{2}/\sigma _{U^{\ast }}^{2}}}\right) \right]
\theta _{0j} \\
& ={}\mathrm{E}\left[ \Phi \left( 
\frac{\theta _{01}\pi _{01}^{\prime }Z_{i}+\left( \theta _{01}\pi
_{02}+\theta _{02}\right) ^{\prime }W_{i}}{\sqrt{2\sigma _{U^{\ast
}}^{2}-\sigma _{U}^{2}+\theta _{01}^{2}\sigma _{V}^{2}}}\right) \right]
\theta _{0j},
\end{align*}%
where the penultimate equality comes from taking expectation with respect to 
$V_{i}^{\ast }$, and the last equality follows from equation~(\ref{eq:sigma2
w and wo stars}).\hfill $\blacksquare $

\bigskip

\noindent \textbf{Derivation of equation~(\ref{eq:true APE Probit})} is
similar:

\begin{align*}
& APE_{j}^{\text{Pr}}\left( \sigma _{U^{\ast }}^{2}\right) =\mathrm{E}\left[
\phi \left( \frac{\theta _{0}^{\prime }H_{i}^{\ast }}{\sigma _{U^{\ast }}}%
\right) \frac{\theta _{0j}}{\sigma _{U^{\ast }}}\right] =\mathrm{E}\left[
\phi \left( \frac{\theta _{01}\pi _{01}^{\prime }Z_{i}+\left( \theta
_{01}\pi _{02}+\theta _{02}\right) ^{\prime }W_{i}}{\sigma _{U^{\ast }}}+%
\frac{\theta _{01}V_{i}^{\ast }}{\sigma _{U^{\ast }}}\right) \right] \frac{%
\theta _{0j}}{\sigma _{U^{\ast }}} \\
& =\frac{1}{\sqrt{1+\theta _{01}^{2}\sigma _{V^{\ast }}^{2}/\sigma _{U^{\ast
}}^{2}}}\mathrm{E}\left[ \phi \left( \frac{\theta _{01}\pi _{01}^{\prime
}Z_{i}+\left( \theta _{01}\pi _{02}+\theta _{02}\right) ^{\prime }W_{i}}{%
\sigma _{U^{\ast }}\sqrt{1+\theta _{01}^{2}\sigma _{V^{\ast }}^{2}/\sigma
_{U^{\ast }}^{2}}}\right) \right] \frac{\theta _{0j}}{\sigma _{U^{\ast }}} \\
& ={}\mathrm{E}\left[ \phi
\left( \frac{\theta _{01}\pi _{01}^{\prime }Z_{i}+\left( \theta _{01}\pi
_{02}+\theta _{02}\right) ^{\prime }W_{i}}{\sqrt{2\sigma _{U^{\ast
}}^{2}-\sigma _{U}^{2}+\theta _{01}^{2}\sigma _{V}^{2}}}\right) \right] 
\frac{\theta _{0j}}{\sqrt{2\sigma _{U^{\ast }}^{2}-\sigma _{U}^{2}+\theta
_{01}^{2}\sigma _{V}^{2}}}.\qquad \qquad \qquad \qquad \qquad \qquad
\blacksquare
\end{align*}

\end{document}
